\chapter{Campo en el plano focal en un sistema estigmático simple refractivo}
\label{ch:focal}

En la sección presente, se calcula el campo en el plano focal en un sistema estigmático simple, formado por una sola superficie dióptrica. Asumimos que incide un campo esférico desde uno de los puntos estigmáticos $A$ y calcularemos el campo refractado en un plano perpendicular al eje óptico que contenga al otro punto estigmático $A'$.

% ===========================================================================
\section{Caso General}
% ===========================================================================

En esta primera parte no especificaremos aún el campo incidente como aquel emitido desde uno de los puntos estigmáticos. El objetivo es obtener una expresión general para el campo transmitido y propagado, dado un campo inicial cualquiera sobre el plano de referencia $P$. La especialización al caso en que la fuente se encuentra en $A$ se realizará posteriormente.


% ---------------------------------------------------------------------------
\subsection{Campo eléctrico en el plano focal}
% ---------------------------------------------------------------------------

El problema que se resolverá en la presente subsección es el de determinar el campo eléctrico en el plano focal dado un campo inicial cualquiera, el cual posteriormente especificaremos como el campo incidente correspondiente a una fuente puntual centrada en uno de los puntos estigmáticos, pero en principio los cálculos aquí presentados son de una naturaleza general.

\textit{Denotaremos como $\vb E$ en general al campo incidente, y a $\vb E'$ el campo tras la refracción.}

Sea $\vb{E}_p(\vb{r})$ el campo eléctrico en el plano $P$. Calcúlese el campo en $P'$ a una distancia $d$ del plano $P$.

Transfiérase $\vb{E}_p$ a la superficie $\Sigma$ adicionando la fase que acumula el campo en su trayecto. Considerar que el campo describe una línea recta en su recorrido de $P$ a $Q$ y asumir que dicho recorrido es aproximadamente horizontal, justificando esto porque consideraremos como primera aproximación un ángulo $\alpha$ moderado.

Con tales consideraciones, el campo $\vb{E}_{\Sigma}$ sobre $\Sigma$ estará dado por

\begin{equation}
    \vb{E}_{\Sigma}(\vb{r})
    =
    e^{ikz(\vb{r})} \vb{E}_p(\vb{r}).
    \label{eq:Ep_from_E_Sigma}
\end{equation}

Equivalentemente,

\begin{equation}
    \vb{E}_{\Sigma}(r,\phi,z(r))
    =
    e^{ikz(r)} \vb{E}_p(r,\phi,z=0).
\end{equation}

Así, calculamos el campo transmitido sobre la interfaz $\Sigma$:

\begin{equation}
    \vb{E}'_{\Sigma}(\vb{q})
    =
    t_p (\vb{E}_{\Sigma} \cdot \uvec{p}) \uvec{p}'
    +
    t_s (\vb{E}_{\Sigma} \cdot \uvec{s}) \uvec{s}.
\label{eq:E_Sigma_Prime_in_terms_of_E_Sigma}
\end{equation}

Descomponiendo el campo en la interfaz en sus componentes cilíndricas,

\begin{equation}
    \vb{E}_{\Sigma}
    =
    E_r \uvec{r}
    +
    E_{\phi} \uvec{\phi}
    +
    E_z \uvec{z}.
\end{equation}

Calculamos las proyecciones sobre los vectores de la base de Fresnel $\uvec p$ y $\uvec \phi$, donde $\uvec{s}=\uvec{\phi}$:

\begin{equation}
    \vb{E}_{\Sigma} \cdot \uvec{p}
    =
    \frac{1}{\sqrt{\rho^2 - 2zz_0 + z_0^2}}
    \big\{
    (z(\rho) - z_0) E_r
    -
    r(\rho) E_z
    \big\},
\end{equation}

\begin{equation}
    \vb{E}_{\Sigma} \cdot \uvec{\phi}
    =
    E_{\phi}.
\end{equation}

Sustituyendo en la expresión del campo transmitido,

\begin{equation}
    \vb{E}'_{\Sigma}
    =
    t_p
    \frac{
    \big\{ (z(\rho) - z_0) E_r - r(\rho) E_z \big\}
    }
    {
    \sqrt{\rho^2 - 2zz_0 + z_0^2}
    \sqrt{\rho^2 - 2zz_i + z_i^2}
    }
    \big\{
    (z(\rho) - z_i) \uvec{r}
    -
    r(\rho) \uvec{z}
    \big\}
    +
    t_s E_{\phi} \uvec{\phi}.
\end{equation}

Transfiérase el campo de vuelta a $P$ para hacer uso de los métodos de propagación; llamamos a este ``campo virtual transmitido'' en $P$ como $\vb{E}_P'$:

\begin{equation}
    \vb{E}_p'
    =
    e^{-inkz(r)} \vb{E}_{\Sigma}'.
\end{equation}

Reemplazando según \eqref{eq:E_Sigma_Prime_in_terms_of_E_Sigma},

\begin{align}
    \vb{E}_p'
    &=
    e^{-inkz(r)}
    \big\{
    t_p (\vb{E}_{\Sigma} \cdot \uvec{p}) \uvec{p}'
    +
    t_s (\vb{E}_{\Sigma} \cdot \uvec{\phi}) \uvec{\phi}
    \big\}
    \nonumber
    \\
    &=
    e^{-inkz(r)}
    \big\{
    t_p (e^{ikz(r)} \vb{E}_p \cdot \uvec{p}) \uvec{p}'
    +
    t_s (e^{ikz(r)} \vb{E}_p \cdot \uvec{\phi}) \uvec{\phi}
    \big\}.
\end{align}

Por tanto,

\begin{equation}
    \vb{E}_p'
    =
    e^{ik(1-n)z(r)}
    \big\{
    t_p (\vb{E}_p \cdot \uvec{p}) \uvec{p}'
    +
    t_s (\vb{E}_p \cdot \uvec{\phi}) \uvec{\phi}
    \big\}.
\end{equation}

Definimos el campo inicial en el plano $z=0$:

\begin{equation}
    \vb{E}(\vb r, z=0)
    :=
    \vb{E}_p(\vb r).
\end{equation}

Y lo propagamos hasta el plano focal $z=f$:

\begin{equation}
    \vb{E}'(r, z=f)
    =
    \mathscr{R}_n [f] \big\{ \vb{E}_p'(r) \big\},
\end{equation}

donde $\mathscr R_n[f]$ denota el operador de propagación que toma un campo en un plano dado y da el campo en otro plano paralelo y a una distancia $f$ al primero en un medio de índice de refracción $n$.

Tomando como válida en nuestro estudio la aproximación de Fresnel, que hemos usado implícitamente cuando transferimos los campos del plano al dioptrio y viceversa, tenemos que el operador de propagación está dado por

\begin{equation}
    \mathscr{R}_n [f] \{ U(x,y) \}
    =
    n\frac{e^{inkf}}{i \lambda f}
    \iint_{-\infty}^{\infty}
    U(x',y')
    \exp\!\left[
        \frac{ink}{2f}
        \lVert \vb{r}-\vb{r}' \rVert^2
    \right]
    d\vb{r}'.
\end{equation}

Por lo tanto, el campo eléctrico en el plano focal está dado por

\begin{equation}
    \vb{E}'(r, z=f)
    =
    n\frac{e^{inkf}}{i \lambda f}
    \iint_{-\infty}^{\infty}
    \vb{E}_p'(r')
    \exp\!\left[
        \frac{ink}{2f}
        \lVert \vb{r}-\vb{r}' \rVert^2
    \right]
    d\vb{r}'.
\end{equation}

O equivalentemente,

\begin{equation}
    \vb{E}'(r, z)
    =
    n\frac{e^{inkz}}{i \lambda z}
    \iint_{-\infty}^{\infty}
    \exp\!\left[
        \frac{ink}{2z}
        \lVert \vb{r}-\vb{r}' \rVert^2
        +
        ik(1-n)z(r')
    \right]
    \Big\{
        t_p (\vb{E}_P \cdot \uvec{p}) \uvec{p}'
        +
        t_s (\vb{E}_P \cdot \uvec{\phi}) \uvec{\phi}
    \Big\}
    dx'dy'.
\label{eq:E_field_at_z_3D}
\end{equation}

Esta expresión nos da el campo en el plano paralelo a $P$ a una distancia $z$, dado el campo inicial en $P$.

Este resultado puede ser también expresado en términos de una transformada de Fourier expandiendo el kernel de Fresnel:

\begin{equation}
\mathbf{E}'(\mathbf r,z)
=
n\frac{e^{inkz}}{i\lambda z}\,
\exp\!\left[
\frac{ink}{2z}\|\mathbf r\|^2
\right]
\,
\mathscr F\{\mathbf P\}
\left(
\frac{n\mathbf r}{\lambda z}
\right).
\label{E_in_observation_plane_as_FT_3D}
\end{equation}

Donde,

\begin{equation}
\mathbf P(\|\mathbf r'\|)
=
\exp\!\left[
\frac{ink}{2z}\|\mathbf r'\|^2
+
ik(1-n)z(\|\mathbf r'\|)
\right]
\left\{
t_p (\mathbf{E}_P\cdot \hat{\boldsymbol p}) \hat{\boldsymbol p}'
+
t_s (\mathbf{E}_P\cdot \hat{\boldsymbol\phi}) \hat{\boldsymbol\phi}
\right\}.
\label{eq:P_definition}
\end{equation}

Defínase un sistema coordenado esférico centrado en el foco del sistema, que es el punto $A'$; llámense por $\uvec R_{A'}$, $\uvec \theta_{A'}$, $\uvec \phi$ los vectores curvilíneos asociados al sistema.

Los estados de polarización $\uvec p'$ y $\uvec s$ son la polarización definida por los vectores base del sistema esférico centrado en $A'$ tangentes a la esfera, esto es $\uvec \theta_{A'}$ y $\uvec \phi$ respectivamente, como puede ser comprobado con facilidad desde las expresiones \eqref{eq:p_prime_mode_ovoid} y \eqref{eq:s_mode_ovoid}.

De modo que en el sistema esférico centrado en $A'$ la base de polarización queda escrita como

\begin{align}
\uvec{s}
&=
\uvec{\phi}
=
-\sin\phi\,\uvec{x}
+
\cos\phi\,\uvec{y},
\\
\uvec{p}'
&=
\uvec{\theta}_{A'}
=
\cos\theta_{A'}
(\cos\phi\,\uvec{x}+\sin\phi\,\uvec{y})
-
\sin\theta_{A'}\,\uvec{z}.
\label{s_p_in_A_prime_system}
\end{align}

Pasamos pues, sin pérdida de generalidad, al problema transversal; denotamos $\vb E'_\perp$ como la parte transversal de $\vb E'$.

De \eqref{eq:E_Sigma_Prime_in_terms_of_E_Sigma} se sigue que

\begin{equation}
\mathbf{E}'_\perp(\mathbf r,z)
=
n\frac{e^{inkz}}{i\lambda z}\,
\exp\!\left[
\frac{ink}{2z}\|\mathbf r\|^2
\right]
\,
\mathscr F\{\mathbf P_\perp\}
\left(
\frac{n\mathbf r}{\lambda z}
\right).
\label{eq:E_field_as_FT}
\end{equation}

donde

\begin{equation}
\mathbf P_\perp(\|\mathbf r'\|)
=
\exp\!\left[
\frac{ink}{2z}\|\mathbf r'\|^2
+
ik(1-n)z(\|\mathbf r'\|)
\right]
\left\{
t_p E_{P,p} \uvec p'_\perp
+
t_s E_{P,\phi} \uvec \phi
\right\}.
\end{equation}

Aquí $\uvec p'_\perp$ y $\uvec s$ denotan las proyecciones transversales de las polarizaciones $\uvec p'$ y $\uvec s$ respectivamente.

% ===========================================================================
\section{Caso particular: sistema estigmático cartesiano con fuente puntual en $A$}
% ===========================================================================

Para analizar nuestro sistema de interés ---un sistema estigmático formado de un único dioptrio---, la función $z(r)$ debe corresponder a la sagita de una superficie cartesiana, y el campo inicial debe corresponder a uno que se emita desde el punto estigmático $A$.

% ---------------------------------------------------------------------------
\subsection{Geometría del dioptrio estigmático}
% ---------------------------------------------------------------------------

Partimos primero por escribir las cantidades geométricas relevantes de las superficies cartesianas para nuestro problema.

\subsubsection{Parametrización de la superficie cartesiana}

Las ecuaciones paramétricas de las superficies dióptricas son

\begin{equation}
  z(\rho)
  =
  \frac{(O + T\rho^2)\rho^2}
  {
  1 + S\rho^2 + \sqrt{1 + (2S - O^2G)\rho^2}
  },
  \label{eq:parametrical_ovoid_z_app}
\end{equation}

\begin{equation}
  r(\rho)
  =
  \sqrt{\rho^2 - z^2(\rho)}.
  \label{eq:parametrical_ovoid_r_app}
\end{equation}

\subsubsection{Vectores de incidencia, transmisión y normal}

El vector normal unitario en cada punto de la superficie considerada es

\begin{equation}
  \mathbf{\hat N}(\rho,\varphi)
  =
  \frac{
  r'(\rho)\,\mathbf{\hat z}
  -
  z'(\rho)\,\mathbf{\hat r}(\varphi)
  }
  {
  \sqrt{(r'(\rho))^2+(z'(\rho))^2}
  }.
  \label{eq:N_ovoid_rho_phi_app}
\end{equation}

Los vectores unitarios de incidencia y transmisión del sistema son

\begin{equation}
  \uvec{u}
  =
  \frac{
  (z(\rho) - z_0)\uvec{z}
  +
  r(\rho)\uvec{r}(\varphi)
  }
  {
  \sqrt{\rho^2 - 2z(\rho)z_0 + z_0^2}
  },
  \label{eq:u_for_P_in_optical_axis_app}
\end{equation}

\begin{equation}
  \uvec{u}'
  =
  \frac{
  (z(\rho) - z_i)\uvec{z}
  +
  r(\rho)\uvec{r}(\varphi)
  }
  {
  \sqrt{\rho^2 - 2z(\rho)z_i + z_i^2}
  }.
  \label{eq:u_prime_for_P_in_optical_axis_app}
\end{equation}

\subsubsection{Base local de Fresnel}

De la teoría electromagnética, sabemos que para sistemas axialmente simétricos el campo sobre la interfaz $\Sigma$ que se transmite se calcula como

\begin{equation}
  \vb{E}'_{\Sigma}
  =
  t_p (\vb{E}_\Sigma \cdot \uvec{p}) \uvec{p}'
  +
  t_s (\vb{E}_\Sigma \cdot \uvec s) \uvec{s},
  \label{eq:Transmited_E_field_in_Sigma_app}
\end{equation}

donde $\uvec s$, $\uvec p$ y $\uvec p'$ se definen como

\begin{equation}
    \uvec s
    =
    \frac{\uvec N \times \uvec u}{\lvert  \uvec N \times \uvec u \rvert},
    \label{eq:s_hat_def_app}
\end{equation}

\begin{equation}
    \uvec p
    =
    \uvec s \times \uvec u,
    \label{eq:p_hat_def_app}
\end{equation}

\begin{equation}
    \uvec p'
    =
    \uvec s \times \uvec u'.
    \label{eq:p_prime_hat_def}
\end{equation}

Reemplazando por las ecuaciones anteriores, se obtiene

\begin{equation}
  \uvec p
  =
  \frac{1}{\sqrt{\rho^2 - 2zz_0 + z_0^2}}
  \left\{
  (z(\rho) - z_0)\uvec{r}(\varphi)
  -
  r(\rho)\uvec{z}
  \right\},
  \label{eq:p_mode_ovoid}
\end{equation}

\begin{equation}
  \uvec{p}'
  =
  \frac{1}{\sqrt{\rho^2 - 2zz_i + z_i^2}}
  \left\{
  (z(\rho) - z_i)\uvec{r}(\varphi)
  -
  r(\rho)\uvec{z}
  \right\},
  \label{eq:p_prime_mode_ovoid}
\end{equation}

\begin{equation}
  \uvec{s}
  =
  \uvec{\phi}(\varphi).
  \label{eq:s_mode_ovoid}
\end{equation}

\subsubsection{Cosenos de incidencia y transmisión}

Los coeficientes de Fresnel de transmisión se definen en función de los cosenos de los ángulos de incidencia y de refracción como

\begin{equation}
t_s
=
\frac{2 n_0 \cos \theta_i}
{n_0 \cos \theta_i + n_i \cos \theta_t},
\label{eq:ts_app}
\end{equation}

\begin{equation}
t_p
=
\frac{2 n_0 \cos \theta_i}
{n_i \cos \theta_i + n_0 \cos \theta_t}.
\label{eq:tp_app}
\end{equation}

Los cosenos $\cos\theta_i$ y $\cos\theta_t$ los obtenemos a partir de las siguientes expresiones:

\begin{align}
  \cos\theta_i &= -\uvec{N} \cdot \uvec{u},
  \\
  \cos\theta_t &= -\uvec{N} \cdot \uvec{u}'.
\end{align}

Explícitamente,

\begin{equation}
  \uvec N(\rho,\varphi)
  =
  \frac{
  r'(\rho)\,\uvec z
  -
  z'(\rho)\,\uvec r(\varphi)
  }
  {
  \sqrt{(r'(\rho))^2+(z'(\rho))^2}
  }.
\label{eq:N_ovoid_rho_phi_explicit}
\end{equation}

Por tanto,

\begin{align}
\cos\theta_i(\rho)
&=
-\,\uvec N(\rho,\varphi)\cdot\uvec u(\rho,\varphi)
\label{eq:cos_theta_i_def_app}
\\
&=
-\;
\frac{
\big(r'(\rho)\uvec z-z'(\rho)\uvec r(\varphi)\big)
\cdot
\big((z(\rho)-z_0)\uvec z+r(\rho)\uvec r(\varphi)\big)
}
{
\sqrt{(r'(\rho))^2+(z'(\rho))^2}\;
\sqrt{\rho^2-2z(\rho)z_0+z_0^2}
}.
\label{eq:cos_theta_i_closed_app}
\end{align}

De donde se obtiene

\begin{equation}
  \cos\theta_i(\rho)
  =
  \frac{
  z'(\rho)\,r(\rho)
  -
  r'(\rho)\,\big(z(\rho)-z_0\big)
  }
  {
  \sqrt{(r'(\rho))^2+(z'(\rho))^2}\;
  \sqrt{\rho^2-2z(\rho)z_0+z_0^2}
  }.
  \label{eq:cos_theta_i_ovoid_z_and_r}
\end{equation}

Usando $r^2(\rho)=\rho^2-z^2(\rho)$, esta expresión puede escribirse como

\begin{equation}
\cos\theta_i(\rho)
=
\frac{
z'(\rho)\,\big(\rho^2-z(\rho)z_0\big)
-
\rho\big(z(\rho)-z_0\big)
}
{
\sqrt{\rho^2\big(1+z'(\rho)^2\big)-2\rho z(\rho)z'(\rho)}
\;
\sqrt{\rho^2-2z(\rho)z_0+z_0^2}
}.
\label{eq:cos_theta_i_ovoid}
\end{equation}

Luego, $\cos\theta_t$ se obtiene simplemente reemplazando $z_0$ por $z_i$:

\begin{equation}
\cos\theta_t(\rho)
=
\frac{
z'(\rho)\,\big(\rho^2-z(\rho)z_i\big)
-
\rho\big(z(\rho)-z_i\big)
}
{
\sqrt{\rho^2\big(1+z'(\rho)^2\big)-2\rho z(\rho)z'(\rho)}
\;
\sqrt{\rho^2-2z(\rho)z_i+z_i^2}
}.
\label{eq:cos_theta_t_ovoid}
\end{equation}

Así pues, los cosenos de incidencia y transmisión son

\begin{equation}
\left\{
\begin{aligned}
\cos\theta_i
&=
\frac{
z'\,\big(\rho^2 - z\,z_0\big) - \rho\big(z - z_0\big)
}
{
\sqrt{\rho^2\big(1+z'^2\big) - 2\rho z z'}\;
\sqrt{\rho^2 - 2 z z_0 + z_0^2}
},
\\
\cos\theta_t
&=
\frac{
z'\,\big(\rho^2 - z\,z_i\big) - \rho\big(z - z_i\big)
}
{
\sqrt{\rho^2\big(1+z'^2\big) - 2\rho z z'}\;
\sqrt{\rho^2 - 2 z z_i + z_i^2}
}.
\end{aligned}
\right.
\label{cosines_ovoid_compact}
\end{equation}

Usando las definiciones

\begin{align}
l_0
&=
\sqrt{\rho^2 - 2zz_0 + z_0^2},
\\
l_i
&=
\sqrt{\rho^2 - 2zz_i + z_i^2},
\label{eq:l0_li_definitions}
\end{align}

que son las distancias $\overline{AQ}$ y $\overline{A'Q}$ respectivamente, y definiendo también

\begin{equation}
  s'
  :=
  \frac{ds}{d\rho}
  =
  \sqrt{(z')^2 + (r')^2},
  \label{eq:ds_drho}
\end{equation}

los cosenos de incidencia y transmisión quedan expresados como

\begin{equation}
\begin{aligned}
  \cos\theta_i
  &=
  \frac{1}{s'l_0}
  \big\{
  z'\big(\rho^2 - z z_0\big)
  -
  \rho\big(z - z_0\big)
  \big\},
\\
  \cos\theta_t
  &=
  \frac{1}{s'l_i}
  \big\{
  z'\big(\rho^2 - z z_i\big)
  -
  \rho\big(z - z_i\big)
  \big\}.
\end{aligned}
\label{cosines_ovoid}
\end{equation}

\subsubsection{Coeficientes de Fresnel sobre la superficie}

Reemplazando \eqref{cosines_ovoid_compact} en \eqref{eq:ts_app} y en \eqref{eq:tp_app}, obtenemos los coeficientes de Fresnel para la transmisión para la superficie cartesiana parametrizada por los parámetros $G,O,T,S$:

\begin{equation}
 t_s
 =
 \frac{
 2 \left\{ z'(\rho^2 - zz_0) - \rho(z - z_0) \right\}
 }
 {
 \left\{ z'(\rho^2 - zz_0) - \rho(z - z_0) \right\}
 +
 \eta \frac{l_0}{l_i}
 \left\{ z'(\rho^2 - zz_i) - \rho(z - z_i) \right\}
 },
  \label{eq:t_s_ovoid}
\end{equation}

\begin{equation}
 t_p
 =
 \frac{
 2 \left\{ z'(\rho^2 - zz_0) - \rho(z - z_0) \right\}
 }
 {
 \eta
 \left\{ z'(\rho^2 - zz_0) - \rho(z - z_0) \right\}
 +
 \frac{l_0}{l_i}
 \left\{ z'(\rho^2 - zz_i) - \rho(z - z_i) \right\}
 }.
  \label{eq:t_p_ovoid}
\end{equation}

Con esto, terminamos de proporcionar la forma exacta de las cantidades relevantes necesarias para el cálculo del campo eléctrico en el plano focal en función de la función sagita que define el dioptrio estigmático de interés.

% ---------------------------------------------------------------------------
\subsection{Forma general del campo incidente cuya fuente es $A$}
\label{subsec:campo_incidente_fuente_A}
% ---------------------------------------------------------------------------

El campo incidente es aquel que tiene como fuente el primer punto estigmático, el punto $A$. Asumiendo que el campo es uno esférico con centro en $A$, en general escribimos el campo incidente $\vb E$ en un punto $V$ cualquiera como

\begin{equation}
  \vb E (V)
  =
  \mathbb{E}_0(V)
  \frac{e^{ik\overline{AV}}}{\overline{AV}},
  \label{eq:initial_field}
\end{equation}

que evaluado en $\Sigma$ y usando la definición \eqref{eq:l0_li_definitions},

\begin{equation}
  \vb E_\Sigma (\vb r)
  =
  \mathbb{E}_0(\vb r)
  \frac{e^{ik\ell_0}}{\ell_0},
  \label{eq:initial_field_sigma}
\end{equation}

donde $\ell_0$ se define según \eqref{eq:l0_li_definitions}, y $\mathbb E_0$ es

\begin{equation}
  \mathbb E_0 (\vb r)
  =
  \mathscr E_\theta \uvec \theta_A
  +
  \mathscr E_\phi \uvec \phi_A,
  \label{eq:mathbb_E_0_definition}
\end{equation}

que asegura que se cumpla la condición de transversalidad. De forma equivalente se escribe

\begin{equation}
  \mathbb E_0 (\vb r)
  =
  \mathscr E_\theta \uvec p
  +
  \mathscr E_\phi \uvec s.
  \label{eq:mathbb_E_0_definition_ps}
\end{equation}

% Nótese que $\mathbb E_0 = \ell_0\,\mathbb E^{\mathscr G}$, en términos de la esfera de referencia definida en \eqref{eq:E_G_def}.

De la relación \eqref{eq:Ep_from_E_Sigma},

\begin{equation}
  \vb E_p (\vb r)
  =
  e^{-ikz(r)}
  \mathbb{E}_0(\vb r)
  \frac{e^{ik\ell_0}}{\ell_0}.
  \label{eq:initial_field_in_p}
\end{equation}

De esta manera queda expresado el campo incidente sobre el plano $P$ para una fuente puntual ubicada en $A$. Si se desea evaluarlo en otro punto del espacio cualquiera $H$, basta con reemplazar $\ell_0$ en \eqref{eq:initial_field_in_p} por $\overline{AH}$. De cualquier manera, el campo incidente queda determinado por dos funciones reales, $\mathscr E_\theta$ y $\mathscr E_\phi$.

% ---------------------------------------------------------------------------
\subsection{Campo eléctrico en el plano focal}
% ---------------------------------------------------------------------------

Reemplazando en \eqref{eq:E_field_as_FT}, se tiene que

\begin{equation}
\mathbf{E}'_\perp(\mathbf r,z)
=
n\frac{e^{inkz}}{i\lambda z}\,
\exp\!\left[
\frac{ink}{2z}\|\mathbf r\|^2
\right]
\,
\mathscr F
\left\{
\frac{1}{\ell_0}
e^{i\frac{nk}{2 z}r'^2 - inkz(r') + ik\ell_0}
\left\{
t_p \mathbb E_0 \cdot \uvec p \, \uvec p'_\perp
+
t_s \mathbb E_0 \cdot \uvec s \, \uvec s
\right\}
\right\}
\left(
\frac{n\mathbf r}{\lambda z}
\right).
\end{equation}

Dado que las polarizaciones $\uvec p$ y $\uvec s$ coinciden con los vectores angulares de la base esférica centrada en $A$, $\uvec \theta_A$ y $\uvec \phi$, podemos escribir

\begin{equation}
\mathbf{E}'_\perp(\mathbf r,z)
=
n\frac{e^{inkz}}{i\lambda z}e^{i\frac{nk}{2z}r^2}
\,
\mathscr F
\left\{
\frac{1}{\ell_0}
e^{ik\left( \frac{n}{2 z}r'^2 - nz(r') + \ell_0 \right) }
\left\{
t_p \mathscr E_{\theta} \uvec{r}
+
t_s \mathscr E_{\phi} \uvec \phi
\right\}
\right\}
\left(
\frac{n\mathbf r}{\lambda z}
\right).
\label{eq:E_perp_for_ovoid_in_prog}
\end{equation}


Identificamos



\begin{equation}
  \Phi(r)
  =
  \frac{n}{2z}r^2
  -
  nz(r)
  +
  \ell_0,
\end{equation}

\begin{equation}
  \mathbb E^{\mathscr G'}
  =
  \frac{1}{\ell_0}
  \left\{
  t_p \mathscr E_\theta \uvec p'
  +
  t_s \mathscr E_\phi \uvec s
  \right\},
  \label{eq:E_G_prime_def}
\end{equation}

\begin{equation}
  \mathbb E^{\mathscr G}
  =
  \frac{1}{\ell_0}
  \left\{
  \mathscr E_\theta \uvec p
  +
  \mathscr E_\phi \uvec s
  \right\},
  \label{eq:E_G_def}
\end{equation}

donde $\mathbb E^{\mathscr G'}$ es el campo sobre la esfera $\mathscr G'$ centrada en el punto imagen $A'$ y tangente al dioptrio, y $\mathbb E^{\mathscr G}$ el campo sobre la esfera $\mathscr G$ centrada en el punto objeto $A$. Dado que la fase es constante sobre $\mathscr G'$, dicho campo es netamente real.

\begin{figure}[H]
  \centering
  \thesisincludegraphics[width=0.65\textwidth]
    {figures/geometrical_setups/EGs.pdf}
    {figures/fast/EGs.jpg}
  \caption{Esferas de referencia $\mathscr G$ y $\mathscr G'$. $\mathscr G$ está centrada en el punto objeto $A$ y $\mathscr G'$ en el punto imagen $A'$; ambas son tangentes al dioptrio en el vértice. El campo $\mathbb E^{\mathscr G}$ es la amplitud incidente sobre $\mathscr G$, y $\mathbb E^{\mathscr G'}$ es el campo transmitido sobre $\mathscr G'$, sobre el que la fase es constante por la condición estigmática.}
  \label{fig:EGs}
\end{figure}

La ecuación \eqref{eq:E_perp_for_ovoid_in_prog} queda entonces expresada como

\begin{equation}
\mathbf{E}'_\perp(\mathbf r,z)
=
n\frac{e^{inkz}}{i\lambda z}e^{i\frac{nk}{2z}r^2}
\,
\mathscr F
\left\{
e^{ik\Phi(r')} \mathbb E^{\mathscr G'}(\vb r')
\right\}
\left(
\frac{n\mathbf r}{\lambda z}
\right).
\label{eq:E_field_perp_no_phase_replacement}
\end{equation}

Para el término $\ell_0$ dividiendo en la amplitud, lo aproximamos simplemente a $z_0$; para la fase la aproximación ha de ser más cuidadosa. En concreto, partiendo de

\begin{equation}
    \ell_0
    =
    \sqrt{x^2 + y^2 + (z(r) - z_0)^2}
    =
    \sqrt{r^2 + (z(r) - z_0)^2},
\end{equation}

extrayendo el término $(z(r) - z_0)$ para realizar una expansión binomial en el régimen paraxial, $r \ll |z(r)-z_0|$:

\begin{align}
    \ell_0
    &=
    (z(r) - z_0)
    \sqrt{
    1
    +
    \left(
    \frac{r}{z(r) - z_0}
    \right)^2
    }
    \nonumber
    \\
    &=
    (z(r) - z_0)
    \left\{
    1
    +
    \frac{1}{2}
    \left(
    \frac{r}{z(r) - z_0}
    \right)^2
    +
    \dots
    \right\}
    \nonumber
    \\
    &\approx
    (z(r) - z_0)
    +
    \frac{1}{2}
    \frac{r^2}{z(r) - z_0}.
\end{align}

Expandiendo el denominador del segundo término, asumiendo que la sagita es pequeña respecto a la distancia al objeto, $z(r)\ll |z_0|$:

\begin{align}
    \ell_0
    &=
    -z_0
    +
    z(r)
    -
    \frac{r^2}{2z_0\left(1-\frac{z(r)}{z_0}\right)}
    \nonumber
    \\
    &=
    -z_0
    +
    z(r)
    -
    \frac{r^2}{2z_0}
    \left\{
    1
    +
    \frac{z(r)}{z_0}
    +
    \left(\frac{z(r)}{z_0}\right)^2
    +
    \dots
    \right\}
    \nonumber
    \\
    \ell_0
    &=
    -z_0
    +
    z(r)
    -
    \frac{r^2}{2z_0}
    +
    \mathcal{O}(r^4).
\end{align}

Recordando la definición de la cantidad $\Phi$,

\begin{equation}
    \Phi
    =
    \frac{n}{2z}r^2
    +
    \ell_0
    -
    nz(r),
  \label{eq:phase_definition}
\end{equation}

y aplicando la aproximación para $\ell_0$,

\begin{align}
  \Phi
  &=
  \frac{n}{2z}r^2
  +
  \left(
  -z_0
  +
  z(r)
  -
  \frac{r^2}{2z_0}
  \right)
  -
  nz(r)
  +
  \mathscr O(r^4)
  \nonumber
  \\
  &=
  -z_0
  +
  (1-n)z(r)
  +
  \frac{1}{2}
  \left(
  \frac{n}{z}
  -
  \frac{1}{z_0}
  \right)r^2
  +
  \mathscr O(r^4).
\end{align}

Procedemos reemplazando la forma de la sagita $z(r)$ en la aproximación paraxial:

\begin{equation}
    z(r)
    \approx
    \frac{O}{2}r^2
    +
    \mathcal{O}(r^4),
\end{equation}

donde $O$ es uno de los parámetros de la parametrización GOTS de las superficies cartesianas que se desarrolló en \cite{singlet2020}, cuyas definiciones escribimos en el presente trabajo en el apéndice \ref{apx:ovoid_implicit_deduction}, y el desarrollo en series de la sagita, $\ell_0$, entre otras cantidades relevantes, se realiza con más detalle en \ref{apx:ovoids_series}. Asumiendo que el índice inicial es $n_i=n$, el índice de salida es $n_o=1$, y la posición de la imagen estigmática es $z_i=f$, el parámetro $O$ está dado por

\begin{equation}
    O
    =
    \frac{n_i z_0 - n_o z_i}{z_i z_0 (n_i - n_o)}
    =
    \frac{1}{n - 1}
    \left(
    \frac{n}{f}
    -
    \frac{1}{z_0}
    \right),
\end{equation}

de modo que

\begin{equation}
    (1-n)O
    =
    \frac{1}{z_0}
    -
    \frac{n}{f}.
\end{equation}

Reemplazando el término de la sagita en la expresión de la fase total $\Phi$:

\begin{align}
    \Phi
    &=
    -z_0
    +
    \frac{1}{2}
    \left(
    \frac{n}{z}
    -
    \frac{1}{z_0}
    \right)r^2
    +
    (1-n)
    \left(
    \frac{O}{2}r^2
    \right)
    \nonumber
    \\
    &=
    -z_0
    +
    \frac{1}{2}
    \left(
    \frac{n}{z}
    -
    \frac{1}{z_0}
    \right)r^2
    +
    \frac{1}{2}
    \left(
    \frac{1}{z_0}
    -
    \frac{n}{f}
    \right)r^2
    \nonumber
    \\
    &=
    -z_0
    +
    \frac{1}{2}
    \left(
    \frac{n}{z}
    -
    \frac{n}{f}
    \right)r^2
    +
    \mathcal{O}(r^4).
\end{align}

Es evidente a partir de este resultado que, cuando la distancia de observación $z$ coincide con el foco del sistema, $z=f$, el término cuadrático en $r$ de la fase se cancela exactamente. Esto significa que la fase de la pupila se vuelve constante, $\Phi\simeq -z_0$, y podemos hacerla $1$ sin pérdida de generalidad.

Reemplazando la fase de pupila por una fase constante en la ecuación \eqref{eq:E_field_perp_no_phase_replacement}, obtenemos

\begin{equation}
  \vb E'(\vb r, f)
  =
  \frac{n}{i\lambda f}e^{inkf}e^{i\frac{nk}{2f}r^2}
  \mathscr F
  \Big\{
  \mathbb E^{\mathscr G'}(r', \phi')
  \Big\}
  \left(
  \frac{n}{\lambda f}\vb r
  \right).
  \label{eq:E_field_pupil_phase_and_TT}
\end{equation}

Llamamos $\mathscr{D}$ a la esfera centrada en el vértice del dioptrio con radio $f$. Su fase paraxial evaluada en el plano focal es precisamente el factor $e^{i\frac{nk}{2f}r^2}$ presente en \eqref{eq:E_field_pupil_phase_and_TT}. Definimos el campo sobre $\mathscr{D}$ como

\begin{equation}
  \vb E^{\mathscr D}(\vb r)
  \coloneqq
  \frac{n}{i\lambda f}\,e^{inkf}
  \,\mathscr F\!\left\{\mathbb E^{\mathscr G'}\right\}
  \!\left(\frac{n}{\lambda f}\vb r\right),
  \label{eq:E_D_def}
\end{equation}

de modo que \eqref{eq:E_field_pupil_phase_and_TT} se reescribe compactamente como

\begin{equation}
  \vb E'_\perp(\vb r, f)
  =
  e^{i\frac{nk}{2f}r^2}\,
  \vb E^{\mathscr D}(\vb r).
  \label{eq:E_field_as_Debye}
\end{equation}

En síntesis, los resultados centrales del capítulo en términos de $\mathbb E^{\mathscr G'}$ son


\begin{equation}
  \vb E'(\vb r, f)
  =
  \frac{ne^{inkf}}{i\lambda f}e^{i\frac{nk}{2f}r^2}
  \,\mathscr F\!\left\{\mathbb E^{\mathscr G'}\right\}\!\left(\frac{n}{\lambda f}\vb r\right)
  \label{eq:E_field_in_EGp}
\end{equation}




% === BOXED EQ ===
\begin{equation}
\boxed{
  \vb E'^{\mathscr D}(\vb r)
  =
  \frac{n e^{inkf}}{i\lambda f}  \,\mathscr F\!\left\{\mathbb E^{\mathscr G'}\right\}\!\left(\frac{n}{\lambda f}\vb r\right)
}
  \label{eq:E_D_in_EGp}
\end{equation}
% ===============

que, usando $\mathbb E^{\mathscr G'} = \mathbb T \mathbb E_0 / \ell_0$ y aproximando $\ell_0 \approx z_0$ en la amplitud, se leen equivalentemente como


% === BOXED EQ ===

\begin{equation}
\boxed{
  \vb E'(\vb r, f)
  =
  \frac{ne^{inkf}}{i\lambda f z_0}\,e^{i\frac{nk}{2f}r^2}\,\mathscr F\!\left\{\mathbb T\mathbb E_0\right\}\!\left(\frac{n}{\lambda f}\vb r\right)
}
  \label{eq:E_prime_func_of_M_E0}
\end{equation}
% ===============


Que en la Esfera de Fourier $\mathscr D$ de $\mathscr G'$ el campo está dado por



\begin{equation}
  \vb E'^{\mathscr D}(\vb r)
  =
  \frac{ne^{inkf}
}{i\lambda f z_0}  \,\mathscr F\!\left\{\mathbb T\mathbb E_0\right\}\!\left(\frac{n}{\lambda f}\vb r\right)
  \label{eq:E_prime_D_function_of_M_E0}
\end{equation}





El campo eléctrico en el plano focal queda determinado por la transformada de Fourier de $\mathbb E^{\mathscr G'} = \mathbb T \mathbb E^{\mathscr G}$, donde

\begin{equation}
  \mathbb T
  =
  \ket{p'} \bra p\, t_p
  +
  \ket{s} \bra s\, t_s.
  \label{eq:def_M_bb}
\end{equation}

Dado que en el frente de onda inicial la fase es constante, era de esperar que el frente de onda transmitido también tuviese fase constante, pues el mismo camino óptico se recorrió, como se sigue de la condición estigmática; no hay, pues, aberración de fase. Sin embargo, el campo en el plano focal es proporcional a $\mathscr F\{\mathbb T \mathbb E_0\} = \mathscr F\{\mathbb T\} * \mathscr F\{\mathbb E_0\}$, lo que hace que incluso cuando $\mathbb E_0$ es constante —y su transformada de Fourier es en consecuencia una delta de Dirac— la convolución con $\mathscr F\{\mathbb T\}$ reproduce $\mathscr F\{\mathbb T\}$ evaluada en el punto de observación, que no es en general una delta dado que $\mathbb T$ es diferente en cada punto considerado, pues depende del coseno de incidencia. Esto produce un ensanchamiento del punto imagen, además de un acoplamiento entre las polarizaciones, efectos que pueden ser entendidos como una aberración de origen netamente vectorial, presentes incluso en sistemas libres de imperfecciones geométricas. 

Dicho de forma breve, el origen de las aberraciones vectoriales se debe a que los coeficientes de Fresnel transforman el campo de manera no homogénea durante la refracción, y lo hacen de manera distinta para dos componentes "principales", a saber, las correspondientes a las polarizaciones $\uvec s$ y $\uvec p$.

\medskip

\noindent\textit{Relación con la formulación de Richards y Wolf.} El resultado
\eqref{eq:E_prime_func_of_M_E0} tiene la misma estructura de la fórmula
canónica del campo focal derivada por Richards y Wolf~\cite{RichardsWolf1959}
para sistemas aplanáticos idealizados: una transformada de Fourier
bidimensional sobre la esfera de referencia imagen, aplicada al producto de
la función pupila escalar por un operador vectorial que gira los estados de
polarización. La diferencia esencial es que en la formulación de Richards y
Wolf el operador vectorial se postula desde la construcción aplanática
—asume una lente ideal que satisface la condición de los senos de Abbe— y
la amplitud sobre la pupila queda ponderada por el factor apodizador
$\sqrt{\cos\theta}$ característico. En cambio, la función pupila vectorial
$\mathbb E^{\mathscr G'}=\mathbb T\mathbb E_0/\ell_0$ derivada aquí se
construye punto a punto desde la geometría exacta de la superficie
cartesiana y desde la transmisión de Fresnel evaluada sobre ella, sin
suponer ninguna lente ideal: la factorización $1/\ell_0$ absorbe la
divergencia esférica sobre $\mathscr G$, y las variaciones angulares de
$t_p$ y $t_s$ producen tanto la modulación de amplitud como el acoplamiento
vectorial. En el límite en que la superficie cartesiana se aproxima por su
esfera osculatriz y $t_p,t_s$ se toman constantes, la función pupila aquí
obtenida colapsa sobre la de Richards y Wolf; la formulación presente es,
por tanto, una generalización que incorpora la geometría real del sistema
estigmático y la transmisión de Fresnel completa.

\medskip

\noindent\textit{Límite escalar y definición operativa de la aberración vectorial.}
El resultado \eqref{eq:E_prime_func_of_M_E0} recupera la teoría escalar de
la difracción de forma exacta al imponer $t_p=t_s=1$, en cuyo caso
$\mathbb T$ se reduce a la identidad transversal y la función pupila se
convierte en una amplitud escalar isótropa sobre $\mathscr G'$. Este límite
no es un régimen físicamente accesible —los coeficientes $t_p$ y $t_s$
solo coinciden con la unidad simultáneamente en el caso trivial de un mismo
medio a ambos lados de la interfaz— sino un límite matemático de
referencia: la desviación del campo focal completo respecto a él es lo que
llamamos aberración vectorial. Adoptamos esta definición porque el
propósito del trabajo es aislar la contribución de la refracción vectorial,
sin mezclarla con otros efectos residuales (contraste de índice, apodización
paraxial, etc.). La cantidad que la gobierna es la combinación anisotrópica
$t_-=t_p-t_s$; el límite escalar corresponde a $t_-\equiv 0$.

\medskip

\noindent\textit{Propagación de Fresnel y régimen de apertura numérica alta.}
El propagador \eqref{eq:E_prime_func_of_M_E0} corresponde al régimen de
Fresnel en el espacio libre entre la esfera de referencia imagen y el plano
focal (véase la subsección de aproximación de Fresnel del
capítulo~\ref{ch:ondulatoria}). A pesar de que
los ejemplos de aplicación al final de este capítulo incluyen aperturas
elevadas ($\mathrm{NA}=0.94$), la aproximación de Fresnel se aplica sobre un
tramo corto —el que va de $\mathscr G'$ al plano focal contiguo— y sobre el
campo ya modulado por $\mathbb T$; los efectos vectoriales que estudiamos se
localizan en el entorno del máximo focal, donde las diferencias respecto de
un propagador exacto de espectro angular son subordinadas a la modulación
introducida por $t_-$.

Si bien el caso de reflexión no es tratado en detalle en el presente trabajo, los resultados son en esencia los mismos, bastaría con reemplazar los coeficientes de Fresnel de refracción por los de reflexión.












% ---------------------------------------------------------------------------
\section{Reducción del problema a la parte transversal}
% ---------------------------------------------------------------------------

Si bien el resultado \eqref{eq:E_field_pupil_phase_and_TT} sugiere un campo tridimensional, el problema puede restringirse al plano perpendicular al eje óptico mediante la condición de transversalidad, que establece una dependencia entre el campo transversal y longitudinal, permitiendo determinar el segundo conociendo el primero.

Escribiendo la condición de transversalidad como que el campo ha de ser tangente al vector radial unitario con origen en $A'$,

\begin{equation}
\vb{E}\cdot\uvec{R}_{A'}=0,
\label{eq:transversality_condition}
\end{equation}

lo que implica

\begin{equation}
E_r\sin\theta_{A'} + E_z\cos\theta_{A'} = 0,
\end{equation}

y por tanto

\begin{equation}
E_z
=
-\tan\theta_{A'}\,E_r.
\label{eq:Ez_from_Er}
\end{equation}

Esto es equivalente a

\begin{equation}
  E_z
  =
  -\frac{r}{z_i} \,E_r.
\label{eq:Ez_from_Er_rz}
\end{equation}

De esta manera, el campo transversal determina completamente el campo longitudinal, y por ende el campo vectorial $\vb E$ completamente.

La intensidad total en la pantalla está dada por

\begin{equation}
  I
  =
  I_\perp
  +
  I_z,
  \label{eq:Intensity_in_Focal_Plane_decomp}
\end{equation}

donde la contribución transversal $I_\perp$ es

\begin{equation}
    I_\perp(\vb{r}, f)
    =
    \|\vb{E}_\perp(\vb{r}, f)\|^2.
\end{equation}

Si la parte transversal del campo se escribe como

\begin{equation}
  \vb E_\perp
  =
  E_r \uvec r
  +
  E_\phi \uvec \phi ,
\end{equation}

entonces

\begin{equation}
  I_\perp(\vb r, f)
  =
  |E_r(\vb r, f)|^2
  +
  |E_\phi(\vb r, f)|^2 .
\end{equation}

La componente longitudinal se determina por la dependencia entre $E_z$ y la parte radial del campo transversal, como se sigue de la ecuación \eqref{eq:Ez_from_Er}. En particular,

\begin{equation}
  E_z(\vb r, f)
  =
  -\tan\theta\,
  E_r(\vb r, f),
\end{equation}

de modo que

\begin{equation}
  I_z(\vb r, f)
  =
  |E_z(\vb r, f)|^2
  =
  \tan^2\theta\,
  |E_r(\vb r, f)|^2 .
\end{equation}

La intensidad total en la pantalla focal no se obtiene sumando cuadráticamente las intensidades, sino sumando las contribuciones ortogonales del campo eléctrico. Por tanto,

\begin{equation}
  I(\vb r, f)
  =
  I_\perp(\vb r, f)
  +
  I_z(\vb r, f).
  \label{eq:Intensity_in_Focal_Plane}
\end{equation}

Usando las expresiones anteriores, queda

\begin{equation}
  I(\vb r, f)
  =
  |E_r(\vb r, f)|^2
  +
  |E_\phi(\vb r, f)|^2
  +
  \tan^2\theta\,
  |E_r(\vb r, f)|^2 ,
\end{equation}

o equivalentemente,

\begin{equation}
  I(\vb r, f)
  =
  \sec^2\theta\,
  |E_r(\vb r, f)|^2
  +
  |E_\phi(\vb r, f)|^2 .
\end{equation}

Queda evidenciado cómo la intensidad del campo en el plano focal depende únicamente de la parte transversal del campo eléctrico.

Como ejemplo de aplicación de la formulación general a un campo sin ninguna
simetría, consideramos la formación de la imagen de una máscara litográfica de
tipo circuito, con trazas rectas en ambas orientaciones cartesianas. El sistema
formador de imagen es una lente formada por dos superficies cartesianas
conjugadas (vacío $\to$ sílice fundida $\to$ vacío), con conjugación exacta y
aumento unitario, de modo que está libre de aberración
geométrica por construcción. Los parámetros corresponden al régimen de la
litografía ultravioleta profunda: $\lambda=193\,\mathrm{nm}$,
$\mathrm{NA}=0.94$, con trazas de anchura $0.64\lambda$ y paso fino de
$1.23\lambda$, cercano al límite de resolución. La máscara es una transmisión
binaria iluminada coherentemente, y se comparan dos modelos de propagación: el
escalar, con $t_p=t_s=1$, y el vectorial completo, que aplica la matriz de
transmisión $\mathbb T$ en cada dioptrio y reconstruye la componente
longitudinal mediante la condición de transversalidad. En ambos casos se
reporta la intensidad total \eqref{eq:Intensity_in_Focal_Plane}. El pipeline
completo de la simulación —composición de operadores en las dos superficies,
grilla, propagación de Fresnel y reconstrucción longitudinal— se detalla en
el apéndice~\ref{apx:numerical_algorithm}.

Para cuantificar la fidelidad de la imagen se mide el contraste
$C=(\bar I_{\mathrm{pico}}-\bar I_{\mathrm{valle}})/(\bar
I_{\mathrm{pico}}+\bar I_{\mathrm{valle}})$ sobre dos cortes que atraviesan
grupos de trazas paralelas: un corte a través de trazas verticales, cuya
modulación espacial es paralela a $\hat{\mathbf x}$ ($C_V$), y otro a través de
trazas horizontales, con modulación paralela a $\hat{\mathbf y}$ ($C_H$). El
modelo escalar predice una fidelidad esencialmente independiente de la
orientación: $C_V=0.91$ y $C_H=0.98$, cuya pequeña diferencia es puramente
geométrica (el grupo vertical es algo más denso). El modelo vectorial rompe
esta equivalencia, como se observa en la
figura~\ref{fig:lithography-general-non-axisymmetric}: con iluminación lineal
$\hat{\mathbf x}$, las trazas alineadas con la polarización conservan la mayor
parte del contraste ($C_H: 0.98\to 0.85$), mientras que las trazas ortogonales
a ella se degradan mucho más severamente ($C_V: 0.91\to 0.67$). Con iluminación
$\hat{\mathbf y}$ los papeles se intercambian ($C_V=0.81$, $C_H=0.68$): los
patrones que se alinean con la polarización lineal incidente se resuelven
mejor que los ortogonales.

El mecanismo dominante en este régimen de apertura es la componente
longitudinal: para las trazas cuya modulación es paralela a la polarización,
las órdenes difractadas se separan en el plano que contiene al campo eléctrico,
sus campos dejan de ser paralelos y la componente $E_z$ —que a
$\mathrm{NA}=0.94$ transporta cerca de la mitad de la energía de la imagen—
rellena los valles entre trazas. Las trazas ortogonales, en cambio, difractan
en el plano perpendicular al campo y no sufren esta proyección. La mezcla de
polarizaciones inducida por $t_-$ es aquí secundaria: la fracción de energía
del canal cruzado es del orden de $0.1\,\%$.

\begin{figure}[H]
  \centering
  \thesisincludegraphics[width=0.95\textwidth]
    {figures/results_of_simulation/vecdiff_final/litografia/litografia_lineal_imagenes.png}
    {figures/fast/litografia_lineal_imagenes.jpg}
  \caption{Imagen de una máscara litográfica de tipo circuito formada por una
  lente de dos superficies cartesianas conjugadas ($\lambda=193\,\mathrm{nm}$,
  $\mathrm{NA}=0.94$, aumento unitario, trazas de $0.64\lambda$ con paso fino
  de $1.23\lambda$). Arriba: máscara de entrada e intensidad predicha por el
  modelo escalar. Abajo: intensidad total del modelo vectorial para
  iluminación lineal $\hat{\mathbf x}$ (izquierda) y $\hat{\mathbf y}$
  (derecha). En cada panel se indican los contrastes $C_V$ y $C_H$ medidos
  sobre los cortes señalados (rojo: trazas verticales; azul: trazas
  horizontales): las trazas alineadas con la polarización incidente se
  resuelven mejor que las ortogonales.}
  \label{fig:lithography-general-non-axisymmetric}
\end{figure}

Por último, la iluminación circular reparte la pérdida de contraste entre ambas
orientaciones ($C_V=0.73$, $C_H=0.74$, figura~\ref{fig:lithography-circular}):
isotropiza el costo vectorial, pero no lo elimina. Este comportamiento —trazas
nítidas en una orientación y degradadas en la ortogonal, en un sistema libre de
aberración geométrica— es una aberración de origen netamente vectorial,
invisible para el modelo escalar, y es el mismo fenómeno que obliga a la
litografía de alta apertura numérica a controlar la polarización de la
iluminación \cite{flagello1996, ruoff2009}.

\begin{figure}[H]
  \centering
  \thesisincludegraphics[width=0.85\textwidth]
    {figures/results_of_simulation/vecdiff_final/litografia/litografia_circular_imagenes.png}
    {figures/fast/litografia_circular_imagenes.jpg}
  \caption{Imagen del mismo patrón de la
  figura~\ref{fig:lithography-general-non-axisymmetric} con iluminación
  circular: modelo escalar (izquierda) y vectorial con intensidad total
  (derecha). La pérdida de contraste se reparte por igual entre ambas
  orientaciones ($C_V=0.73$, $C_H=0.74$): la iluminación circular isotropiza
  la aberración vectorial, pero no la elimina.}
  \label{fig:lithography-circular}
\end{figure}

% ---------------------------------------------------------------------------
\section{Transformación del estado de polarización tras la refracción}
\label{subsec:polarizacion_refraccion}

La ecuación \eqref{eq:def_M_bb} muestra que el campo sobre la esfera de
referencia imagen $\mathscr{G}'$ se obtiene mediante la acción local de la
matriz de transmisión $\mathbb{T}$ sobre el campo en la esfera incidente
$\mathscr{G}$. Hasta este punto hemos enfatizado el efecto de esta matriz
sobre la amplitud espacial del campo y, en consecuencia, sobre el campo
difractado. Sin embargo, $\mathbb{T}$ también actúa como un operador de Jones
local, por lo que modifica punto a punto el estado de polarización. En esta
sección determinamos el estado de polarización sobre $\mathscr{G}'$
conocido el estado sobre $\mathscr{G}$.

Consideremos el estado de Jones incidente normalizado en cada punto
$(\rho',\phi')$ de la pupila,

\begin{equation}
\hat{\vb J}(\rho',\phi')
=
e^{i\gamma(\rho',\phi')}
\begin{pmatrix}
\cos \vartheta(\rho',\phi') \\
\sin \vartheta(\rho',\phi')\, e^{i\delta(\rho',\phi')}
\end{pmatrix}.
\label{eq:jones_state_pupila}
\end{equation}

Aquí $\gamma$ es una fase global, $\vartheta$ determina la razón de
amplitudes entre las componentes cartesianas del campo, y $\delta$ es la fase
relativa. Reservamos $\psi$ para el ángulo de orientación geométrica de la
elipse de polarización y $\chi$ para su ángulo de elipticidad, cuya relación
con el estado se dará al final de la sección en términos de los parámetros
de Stokes.

Para un sistema axialmente simétrico, la matriz de transmisión en la
representación cartesiana transversal adopta la forma

\begin{equation}
\mathbb{T}_{xy}
=
\frac{t_+(\rho')}{2}\,\mathbb{I}_{\perp}
+
\frac{t_-(\rho')}{2}
\begin{pmatrix}
\cos 2\phi' & \sin 2\phi' \\
\sin 2\phi' & -\cos 2\phi'
\end{pmatrix},
\label{eq:T_xy_jones}
\end{equation}

donde $\mathbb{I}_{\perp}$ es la identidad en el plano transversal y hemos
introducido las combinaciones

\begin{equation}
t_+
=
t_p+t_s,
\qquad
t_-
=
t_p-t_s.
\label{eq:t_plus_t_minus_def}
\end{equation}

La deducción detallada de esta forma de $\mathbb{T}_{xy}$ a partir de
\eqref{eq:def_M_bb} se presenta en la
subsección~\ref{subsec:representacion_cartesiana},
ecuación~\eqref{eq:M_in_cartesian_basis}. La dependencia en el ángulo azimutal
$\phi'$ proviene de que los estados propios locales de transmisión son las
polarizaciones $\hat{\vb p}$ y $\hat{\vb s}$, orientadas radial y
azimutalmente respecto al eje óptico. Esta forma de $\mathbb{T}_{xy}$ requiere
la simetría axial del sistema óptico, pero no requiere que el campo incidente
sea axialmente simétrico.

La estructura de \eqref{eq:T_xy_jones} se expone con claridad al descomponer
la matriz en la base de Pauli,

\begin{equation}
\mathbb{T}_{xy}
=
a_0\,\sigma_0 + a_1\,\sigma_1 + a_2\,\sigma_2 + a_3\,\sigma_3,
\label{eq:T_xy_pauli}
\end{equation}

donde $\sigma_0=\mathbb{I}_{\perp}$ y

\begin{equation*}
\sigma_1
=
\begin{pmatrix}
1 & 0 \\ 0 & -1
\end{pmatrix},
\qquad
\sigma_2
=
\begin{pmatrix}
0 & 1 \\ 1 & 0
\end{pmatrix},
\qquad
\sigma_3
=
\begin{pmatrix}
0 & -i \\ i & 0
\end{pmatrix},
\end{equation*}

con coeficientes

\begin{equation}
a_0=\frac{t_+(\rho')}{2},
\qquad
a_1=\frac{t_-(\rho')}{2}\cos 2\phi',
\qquad
a_2=\frac{t_-(\rho')}{2}\sin 2\phi',
\qquad
a_3=0.
\label{eq:T_xy_pauli_coeffs}
\end{equation}

Esta descomposición corresponde a la manera en que Chipman y colaboradores
proponen caracterizar la \emph{aberración de polarización} de un sistema
óptico \cite{mcguire1994}. La analogía con la teoría de aberraciones de frente
de onda es directa: así como la aberración escalar se define como una función
de fase sobre la esfera de referencia de salida —y no sobre el campo en el
plano focal—, la aberración vectorial se describe mediante la función
matricial $\mathbb{T}_{xy}(\rho',\phi')$ definida sobre $\mathscr{G}'$, cuyos
coeficientes de Pauli $a_j(\rho',\phi')$ desempeñan el papel de los
coeficientes de aberración. En nuestro caso, $a_0$ describe la parte escalar
de la transmisión —una apodización idéntica para ambas componentes—; el par
$(a_1,a_2)$, de módulo $|t_-(\rho')|/2$ y dependencia azimutal $2\phi'$,
describe una anisotropía lineal local equivalente a un diatenuador orientado
radialmente en cada punto de la pupila; y $a_3=0$ expresa que la refracción en
una superficie isótropa no introduce quiralidad. Sobre esta descomposición se
construyen, por ejemplo, los polinomios de Zernike de orientación de Ruoff y
Totzeck \cite{ruoff2009}, generalización vectorial de los polinomios de
Zernike escalares empleada para expandir sistemáticamente
$\mathbb{T}_{xy}(\rho',\phi')$.

La acción local sobre el estado de polarización puede representarse de manera
compacta mediante la coordenada proyectiva compleja

\begin{equation}
\zeta(\rho',\phi')
=
\frac{\mathscr E_y(\rho',\phi')}{\mathscr E_x(\rho',\phi')}
=
\tan \vartheta(\rho',\phi')\, e^{i\delta(\rho',\phi')},
\label{eq:zeta_def}
\end{equation}

donde $\mathscr E_x$ y $\mathscr E_y$ son las componentes cartesianas
transversales del campo sobre $\mathscr{G}$. Nótese que $\zeta$ es insensible
a la fase global $\gamma$. Después de la transmisión, el estado sobre
$\mathscr{G}'$ queda caracterizado por

\begin{equation}
\zeta_{\mathscr{G}'}(\rho',\phi')
=
\frac{E_y^{\mathscr{G}'}(\rho',\phi')}{E_x^{\mathscr{G}'}(\rho',\phi')}.
\label{eq:zeta_Gp_def}
\end{equation}

Usando la matriz $\mathbb{T}_{xy}$ se obtiene la transformación fraccional

\begin{equation}
\boxed{
\zeta_{\mathscr{G}'}(\rho',\phi')
=
\frac{
t_-(\rho')\sin 2\phi'
+
\left[t_+(\rho') - t_-(\rho')\cos 2\phi'\right]\zeta(\rho',\phi')
}{
t_+(\rho') + t_-(\rho')\cos 2\phi'
+
t_-(\rho')\sin 2\phi'\,\zeta(\rho',\phi')
}.
}
\label{eq:mobius_transformation}
\end{equation}

Esta expresión muestra que la transmisión vectorial induce una transformación
de Möbius sobre el estado local de polarización, la cual es invertible siempre
que $\det \mathbb{T}_{xy}=t_p t_s \neq 0$, condición que se satisface en
transmisión. Los parámetros transmitidos se recuperan como

\begin{equation}
\vartheta_{\mathscr{G}'}(\rho',\phi')
=
\arctan \left|\zeta_{\mathscr{G}'}(\rho',\phi')\right|,
\qquad
\delta_{\mathscr{G}'}(\rho',\phi')
=
\arg \left[\zeta_{\mathscr{G}'}(\rho',\phi')\right].
\label{eq:jones_params_transmitted}
\end{equation}

De forma equivalente, el estado puede representarse mediante el vector de
Stokes normalizado

\begin{equation}
\vb s
=
\frac{1}{1+|\zeta|^2}
\begin{pmatrix}
1-|\zeta|^2 \\
2\operatorname{Re}\zeta \\
2\operatorname{Im}\zeta
\end{pmatrix},
\label{eq:stokes_from_zeta}
\end{equation}

de modo que la transmisión induce el mapa local
$\vb s(\rho',\phi') \mapsto \vb s_{\mathscr{G}'}(\rho',\phi')$, que se obtiene
al evaluar \eqref{eq:stokes_from_zeta} en $\zeta_{\mathscr{G}'}$. En términos
de estas componentes, los ángulos de la elipse de polarización quedan
determinados por

\begin{equation}
\tan 2\psi = \frac{s_2}{s_1},
\qquad
\sin 2\chi = s_3,
\label{eq:psi_chi_from_stokes}
\end{equation}

relaciones válidas tanto para el estado incidente como para el transmitido.

Por tanto, la matriz de transmisión no solo introduce una modulación espacial
de amplitud, sino también una transformación local del estado de polarización.
En el límite escalar $t_p=t_s$ se tiene $t_-=0$, y la
ecuación~\eqref{eq:mobius_transformation} se reduce a
$\zeta_{\mathscr{G}'}=\zeta$: la transmisión no modifica el estado local de
polarización. En cambio, cuando $t_p\neq t_s$, la diferencia entre las
transmisiones de las componentes $p$ y $s$ produce una redistribución espacial
de la polarización sobre $\mathscr{G}'$.

Cabe enfatizar que esta transformación describe el estado de polarización
sobre la esfera de referencia $\mathscr{G}'$, es decir, inmediatamente después
de la refracción. El estado de polarización en el plano focal no está dado por
una transformación local de este tipo: de acuerdo con
\eqref{eq:E_prime_func_of_M_E0}, el campo en el plano focal resulta de la
síntesis de Fourier, que superpone coherentemente las contribuciones de todos
los puntos de la pupila. El estudio del estado de polarización en el plano
focal se desarrolla en la sección siguiente.

% ---------------------------------------------------------------------------
\section{Transformación del estado de polarización en el plano focal}
\label{subsec:polarizacion_plano_focal}

En la sección~\ref{subsec:polarizacion_refraccion} se determinó el estado
de polarización sobre la esfera de referencia $\mathscr{G}'$, inmediatamente
después de la refracción. Para determinar el estado en el plano focal,
separamos el campo focal \eqref{eq:E_prime_func_of_M_E0} en sus dos
componentes cartesianas,

\begin{equation}
\begin{aligned}
E'_x(\vb r,f)
&=
\frac{n e^{inkf}}{i\lambda f z_0}
e^{i\frac{nk}{2f}r^2}
\mathscr F\left\{(\mathbb{T}_{xy}\mathbb E_0)_x\right\}
\left(\frac{n}{\lambda f}\vb r\right),
\\[0.5em]
E'_y(\vb r,f)
&=
\frac{n e^{inkf}}{i\lambda f z_0}
e^{i\frac{nk}{2f}r^2}
\mathscr F\left\{(\mathbb{T}_{xy}\mathbb E_0)_y\right\}
\left(\frac{n}{\lambda f}\vb r\right),
\end{aligned}
\label{eq:Exy_focal_components}
\end{equation}

donde, de acuerdo con la descomposición \eqref{eq:T_xy_jones},

\begin{equation}
\begin{aligned}
(\mathbb{T}_{xy}\mathbb E_0)_x
&=
\frac{1}{2}
\left[
t_+ \mathscr E_x
+
t_-
\left(
\cos 2\phi'\, \mathscr E_x
+
\sin 2\phi'\, \mathscr E_y
\right)
\right],
\\[0.5em]
(\mathbb{T}_{xy}\mathbb E_0)_y
&=
\frac{1}{2}
\left[
t_+ \mathscr E_y
+
t_-
\left(
\sin 2\phi'\, \mathscr E_x
-
\cos 2\phi'\, \mathscr E_y
\right)
\right].
\end{aligned}
\label{eq:TxyE0_components}
\end{equation}

Las funciones \eqref{eq:TxyE0_components} no son otra cosa que las componentes
cartesianas transversales del campo refractado sobre $\mathscr{G}'$, salvo el
factor común $1/\ell_0 \approx 1/z_0$, pues
$\mathbb E^{\mathscr{G}'}=\mathbb{T}\,\mathbb E_0/\ell_0$. En particular, su
cociente reproduce el estado local tras la refracción,

\begin{equation}
\frac{
(\mathbb{T}_{xy}\mathbb E_0)_y
}{
(\mathbb{T}_{xy}\mathbb E_0)_x
}
=
\frac{
E_y^{\mathscr{G}'}
}{
E_x^{\mathscr{G}'}
}
=
\zeta_{\mathscr{G}'}(\rho',\phi'),
\label{eq:zeta_Gp_ratio}
\end{equation}

dado por la transformación de Möbius \eqref{eq:mobius_transformation}. La
síntesis de Fourier actúa, por tanto, sobre un campo cuyo estado de
polarización ya fue transformado localmente por la refracción.

Cada componente cartesiana del campo en el plano focal queda así expresada
como una única transformada de Fourier escalar. Ahora bien, la transformada de
Fourier es en general una función compleja, aun cuando la función transformada
tenga fase constante. Escribimos entonces cada síntesis en forma polar,

\begin{equation}
\begin{aligned}
\mathscr F\left\{(\mathbb{T}_{xy}\mathbb E_0)_x\right\}
\left(\frac{n}{\lambda f}\vb r\right)
&=
A_x(\vb r)\,e^{i\Phi_x(\vb r)},
\\[0.5em]
\mathscr F\left\{(\mathbb{T}_{xy}\mathbb E_0)_y\right\}
\left(\frac{n}{\lambda f}\vb r\right)
&=
A_y(\vb r)\,e^{i\Phi_y(\vb r)},
\end{aligned}
\label{eq:fourier_polar_form}
\end{equation}

con $A_x,A_y\geq 0$ y $\Phi_x,\Phi_y$ reales; es decir,

\begin{equation}
A_{x,y}(\vb r)
=
\left|
\mathscr F
\left\{
(\mathbb{T}_{xy}\mathbb E_0)_{x,y}
\right\}
\left(\frac{n}{\lambda f}\vb r\right)
\right|,
\qquad
\Phi_{x,y}(\vb r)
=
\arg
\mathscr F
\left\{
(\mathbb{T}_{xy}\mathbb E_0)_{x,y}
\right\}
\left(\frac{n}{\lambda f}\vb r\right).
\label{eq:Axy_Phixy_def}
\end{equation}

La difracción introduce de este modo, sobre cada componente cartesiana, una
función de fase adicional propia, en general distinta para cada una. El
prefactor común de \eqref{eq:Exy_focal_components} —incluida la fase
cuadrática $e^{i\frac{nk}{2f}r^2}$— tiene el mismo módulo y la misma fase para
ambas componentes, por lo que constituye una fase global y no interviene en el
estado de polarización.

En consecuencia, el estado de polarización en cada punto del plano focal queda
completamente determinado por la razón entre las amplitudes y por la
diferencia de fase entre las dos síntesis:

\begin{equation}
\tan \vartheta_f(\vb r)
=
\frac{A_y(\vb r)}{A_x(\vb r)},
\qquad
\delta_f(\vb r)
=
\Phi_y(\vb r)-\Phi_x(\vb r).
\label{eq:theta_delta_focal}
\end{equation}

Esto es, en términos de las transformadas de Fourier de las componentes,

\begin{equation}
\tan \vartheta_f(\vb r)
=
\frac{
\left|
\mathscr F
\left\{
(\mathbb{T}_{xy}\mathbb E_0)_y
\right\}
\right|
}{
\left|
\mathscr F
\left\{
(\mathbb{T}_{xy}\mathbb E_0)_x
\right\}
\right|
},
\qquad
\delta_f(\vb r)
=
\arg
\mathscr F
\left\{
(\mathbb{T}_{xy}\mathbb E_0)_y
\right\}
-
\arg
\mathscr F
\left\{
(\mathbb{T}_{xy}\mathbb E_0)_x
\right\},
\label{eq:theta_delta_focal_explicit}
\end{equation}

con ambas transformadas evaluadas en $\frac{n}{\lambda f}\vb r$.

En términos de la coordenada proyectiva introducida en \eqref{eq:zeta_def}, el
estado en el plano focal se escribe de forma compacta como

\begin{equation}
\boxed{
\zeta_f(\vb r)
=
\tan \vartheta_f(\vb r)\,e^{i\delta_f(\vb r)}
=
\frac{
\mathscr F
\left\{
(\mathbb{T}_{xy}\mathbb E_0)_y
\right\}
}{
\mathscr F
\left\{
(\mathbb{T}_{xy}\mathbb E_0)_x
\right\}
}
\left(
\frac{n}{\lambda f}\vb r
\right),
}
\label{eq:zeta_focal}
\end{equation}

y los parámetros geométricos de la elipse de polarización —el vector de Stokes
normalizado y los ángulos $\psi$ y $\chi$— se obtienen evaluando
\eqref{eq:stokes_from_zeta} y \eqref{eq:psi_chi_from_stokes} en $\zeta_f$.

La conexión con la transformación del estado de polarización al refractarse se
hace explícita escribiendo, según \eqref{eq:zeta_Gp_ratio},

\begin{equation*}
(\mathbb{T}_{xy}\mathbb E_0)_y
=
\zeta_{\mathscr{G}'}(\rho',\phi')\,
(\mathbb{T}_{xy}\mathbb E_0)_x.
\end{equation*}

Sustituyendo en \eqref{eq:zeta_focal}, junto con la forma explícita
\eqref{eq:TxyE0_components} de $(\mathbb{T}_{xy}\mathbb E_0)_x$ —el factor
común $1/2$ se cancela—,

\begin{equation}
\zeta_f(\vb r)
=
\frac{
\mathscr F
\left\{
\zeta_{\mathscr{G}'}
\left[
t_+\mathscr E_x
+
t_-
\left(
\cos 2\phi'\, \mathscr E_x
+
\sin 2\phi'\, \mathscr E_y
\right)
\right]
\right\}
}{
\mathscr F
\left\{
t_+\mathscr E_x
+
t_-
\left(
\cos 2\phi'\, \mathscr E_x
+
\sin 2\phi'\, \mathscr E_y
\right)
\right\}
},
\label{eq:zeta_focal_mobius}
\end{equation}

con ambas transformadas evaluadas en $\frac{n}{\lambda f}\vb r$.

El estado de polarización en cada punto del plano focal es, pues, un promedio
coherente de los estados locales $\zeta_{\mathscr{G}'}(\rho',\phi')$
producidos por la refracción —dados por la transformación de Möbius
\eqref{eq:mobius_transformation}—, con pesos complejos determinados por la
amplitud $(\mathbb{T}_{xy}\mathbb E_0)_x$ y por el núcleo de Fourier. La
transformación local del estado al refractarse queda así incorporada
explícitamente en el resultado focal.

Este resultado es completamente general: no requiere simetría azimutal de las
amplitudes ni uniformidad del estado de polarización sobre la pupila. Nótese
la diferencia esencial con la transformación local de la
sección~\ref{subsec:polarizacion_refraccion}: las amplitudes $A_{x,y}$ y
las fases $\Phi_{x,y}$ en un punto de observación son funcionales de la
distribución completa del campo sobre la pupila, de modo que en general no
existe una relación local entre el estado incidente $\zeta(\rho',\phi')$ y el
estado focal $\zeta_f(\vb r)$.

La lectura física de \eqref{eq:theta_delta_focal} es directa. Aunque el campo
sobre la pupila tenga fase constante —como garantiza la condición
estigmática—, las dos síntesis de Fourier adquieren fases de difracción
$\Phi_x$ y $\Phi_y$ que son funciones distintas del punto de observación. La
diferencia $\delta_f=\Phi_y-\Phi_x$ actúa entonces como un desfasador efectivo
distribuido sobre el plano focal: allí donde toma valores distintos de $0$ y
$\pi$ genera elipticidad, aun para iluminación incidente linealmente
polarizada. La razón $A_y/A_x$, por su parte, redistribuye el peso entre las
componentes y con ello la orientación de la elipse.

Como caso ilustrativo, bajo las simetrías azimutales que se estudiarán en la
subsección~\ref{subsec:representacion_cartesiana}, las síntesis se reducen a
las transformadas de Hankel de \eqref{eq:Ex_Ey_Hankel_cartesian}. Si además el
estado incidente es lineal, las componentes $\mathscr E_x$ y $\mathscr E_y$
pueden tomarse reales, y ambas síntesis resultan funciones reales: las fases
$\Phi_x$ y $\Phi_y$ solo toman los valores $0$ y $\pi$, con saltos de $\pi$ en
los ceros anulares de las transformadas. En tal caso $\delta_f\in\{0,\pm\pi\}$
y el estado permanece lineal en todo punto del plano focal, manifestándose la
aberración vectorial como una rotación local de la orientación, como se
observará en la figura~\ref{fig:lineal-plano-focal-mapa-polarizacion}. Fuera de
estas condiciones de simetría las fases varían de manera continua y la
difracción vectorial genera, en general, elipticidad espacialmente variable.

% ---------------------------------------------------------------------------
\section{Resultados especiales en las distintas bases de polarización}
% ---------------------------------------------------------------------------


Obtenida la expresión matemática para el campo transmitido en el plano focal en \eqref{eq:E_field_pupil_phase_and_TT} en función del campo incidente, procedemos a escribir el resultado en distintas representaciones. Estudiaremos la representación circular, la cilíndrica y, por último, la cartesiana; por representación en una base dada nos referimos a que tanto el campo transmitido como el incidente se indican en dicha base.

No será esto una cuestión únicamente de representación, sino que se estudiará qué sucede cuando las componentes de la representación correspondiente exhiben simetría azimutal,
que no ha de confundirse con simetría azimutal absoluta, pues un campo cuya componente radial sea azimutalmente simétrica no significa que lo sea en el sentido absoluto (que el campo eléctrico sea realmente el mismo en magnitud y dirección, independiente del valor de la coordenada $\phi$), ya que por componente azimutalmente simétrica entendemos la función escalar que multiplica al campo vectorial radial $\vb r(\phi)$, el cual no es azimutalmente simétrico.


Para proceder con los resultados en cada representación, se parte de la ecuación \eqref{eq:E_field_pupil_phase_and_TT} junto con la definición \eqref{eq:E_G_prime_def}, y escribiremos tanto $\mathscr E_\theta$, $\mathscr E_\phi$ como los vectores unitarios $\uvec s$  y $\uvec p'$ en la representación deseada.

En la presente sección tendremos en cuenta las siguientes definiciones.

Definimos la transformada de Hankel de orden $m$ como

\begin{equation}
\mathscr H_m[f](q)
=
\int_0^\infty
f(r')J_m(qr')r'\,dr',
\label{eq:def_Hankel_Transform_m_order}
\end{equation}

donde $q$ es la frecuencia radial conjugada a $r'$ en la transformada de Fourier bidimensional.

También usaremos las combinaciones $t_\pm = t_p \pm t_s$ introducidas en
\eqref{eq:t_plus_t_minus_def}.

Para el caso de una pupila circular de radio $a$, definimos la función pupila

\begin{equation}
A_a(r')
=
\begin{cases}
1, & r' \leq a, \\
0, & r' > a.
\end{cases}
\label{eq:def_pupil_function}
\end{equation}


Procedemos a dar los resultados en las distintas representaciones.


%  ==== REPRESENTACIÓN CIRCULAR ====

% ···········································································
\subsection{Campo eléctrico en el plano focal en la representación circular}
% ···········································································


%\eqref{eq:E_field_pupil_phase_and_TT}




Definiendo los vectores unitarios de polarización circular izquierda y polarización circular derecha, que son las bases de la representación circular como



\begin{equation}
\begin{aligned}
  \ket L
  &=
  \frac{1}{\sqrt 2}
  \left\{
  \uvec x + i \uvec y
  \right\},
  \\
  \ket R
  &=
  \frac{1}{\sqrt 2}
  \left\{
  \uvec x - i \uvec y
  \right\}.
  \label{eq:def_of_circular_basis}
\end{aligned}
\end{equation}



Por lo tanto, los vectores $\uvec s_\perp$ y $\uvec p'_\perp$ en la base circular quedan dados según

\begin{align}
\uvec{s}_\perp
&=
\uvec{\phi}
=
-\sin\phi\,\uvec{x}
+
\cos\phi\,\uvec{y},
\\
\uvec{p}'_\perp
&=
\cos\phi\,\uvec{x}
+
\sin\phi\,\uvec{y}.
\end{align}


Las proyecciones perpendiculares de la base de Fresnel se expresan como

\begin{align}
\uvec{s}_\perp
&=
\frac{i}{\sqrt{2}}
\left(
e^{i\phi}\ket{R}
-
e^{-i\phi}\ket{L}
\right),
\\
\uvec{p}'_\perp
&=
\frac{1}{\sqrt{2}}
\left(
e^{-i\phi}\ket{L}
+
e^{i\phi}\ket{R}
\right).
\end{align}


Reemplazando en \eqref{eq:E_G_prime_def},

\begin{equation}
  \ell_0 \mathbb{E}^{\mathscr G'}_\perp
=
\frac{1}{\sqrt{2}}
\Big\{
(t_p \mathscr E_\theta - i t_s \mathscr E_\phi)e^{-i\phi}\ket{L}
+
(t_p \mathscr E_\theta + i t_s \mathscr E_\phi)e^{i\phi}\ket{R}
\Big\}.
\end{equation}

Para que la integral quede completamente formulada en términos de la base circular, escribimos las componentes del campo incidente en dicha representación:

\begin{align}
\mathscr E_\theta
=
E_r
&=
\frac{1}{\sqrt{2}}
\left(
e^{i\phi}\mathscr E_L
+
e^{-i\phi}\mathscr E_R
\right),
\label{eq:Ep_in_LR_basis}
\\
\mathscr E_\phi
&=
\frac{i}{\sqrt{2}}
\left(
e^{i\phi}\mathscr E_L
-
e^{-i\phi}\mathscr E_R
\right).
\label{eq:Ephi_in_LR_basis}
\end{align}

Por lo tanto,

\begin{equation}
t_p\mathscr E_\theta+i t_s\mathscr E_\phi
=
\frac{1}{\sqrt2}
\left[
(t_p-t_s)e^{i\phi}\mathscr E_L
+
(t_p+t_s)e^{-i\phi}\mathscr E_R
\right],
\end{equation}

\begin{equation}
t_p\mathscr E_\theta-i t_s\mathscr E_\phi
=
\frac{1}{\sqrt2}
\left[
(t_p+t_s)e^{i\phi}\mathscr E_L
+
(t_p-t_s)e^{-i\phi}\mathscr E_R
\right].
\end{equation}

Definimos, en favor de la simple escritura

\begin{equation}
  \vb E_t
  \equiv
  \ell_0 \mathbb E^{\mathscr G'}_\perp.
\end{equation}

De modo que,

\begin{equation}
\vb E_t
=
\frac{1}{2}
\Big\{
\big[(t_p+t_s)\mathscr E_L+(t_p-t_s)e^{-2i\phi}\mathscr E_R\big]\ket L
+
\big[(t_p-t_s)e^{2i\phi}\mathscr E_L+(t_p+t_s)\mathscr E_R\big]\ket R
\Big\}.
\label{eq:T_perp_circular_basis}
\end{equation}

Por lo tanto,

\begin{equation}
\begin{aligned}
T_L
&=
\frac{1}{2}
\big[
(t_p+t_s)\mathscr E_L
+
(t_p-t_s)e^{-2i\phi}\mathscr E_R
\big],
\\
T_R
&=
\frac{1}{2}
\big[
(t_p-t_s)e^{2i\phi}\mathscr E_L
+
(t_p+t_s)\mathscr E_R
\big],
\end{aligned}
\label{eq:TL_TR_circular}
\end{equation}

donde por supuesto, 

\begin{equation}
  \vb E_t
  =
  T_L\ket L
  +
  T_R\ket R.
\end{equation}



Es así que podemos escribir 

\begin{equation}
  \vb E_t = \mathbb T \begin{bmatrix} 
                      \mathscr E_L \\ \mathscr E_R
                    \end{bmatrix},
\end{equation}



donde $\mathbb{T}$ para esta representación es

\begin{equation}
\mathbb T
=
\frac{1}{2}
\begin{bmatrix}
t_+
&
t_-\,e^{-2i\phi'}
\\
t_-\,e^{2i\phi'}
&
t_+
\end{bmatrix}_{LR}.
\label{eq:M_in_circular_basis}
\end{equation}


Reemplazando en \eqref{eq:E_prime_func_of_M_E0} obtenemos el campo transmitido en el plano focal,
en la representación circular




\begin{equation}
\boxed{
\begin{aligned}
\mathbf{E}'_\perp(\mathbf r, f)
=
\frac{n e^{inkf}}{2i\lambda f z_0}
e^{i\frac{nk}{2f}r^2}
\Bigg\{
& \mathscr F\!\left[
t_+\mathscr E_L
+
t_-\,e^{-2i\phi'}\mathscr E_R
\right]
\!\left(
\frac{n}{\lambda f}\mathbf r
\right)
\ket{L}
\\
\qquad\qquad
+
& \mathscr F\!\left[
t_-\,e^{2i\phi'}\mathscr E_L
+
t_+\mathscr E_R
\right]
\!\left(
\frac{n}{\lambda f}\mathbf r
\right)
\ket{R}
\Bigg\}.
\end{aligned}
}
\label{eq:E_prime_in_circular_representation}
\end{equation}


Este resultado da de manera general el campo transmitido evaluado en el plano focal, como función de las funciones escalares $\mathscr E_\theta$ y $\mathscr E_\phi$ que dependen de $r$ y $\phi$ y que, de acuerdo con lo discutido en la subsección~\ref{subsec:campo_incidente_fuente_A}, son funciones reales y definen el vector amplitud del campo en la esfera incidente.



\subsubsection{Reducción del problema a componentes axialmente simétricas}

La ecuación \eqref{eq:E_prime_in_circular_representation} puede simplificarse notablemente si asumimos que las componentes $\mathscr E_L$ y $\mathscr E_R$ presentan simetría axial. Además, son de interés en óptica, pues campos escritos de esa manera representan haces vectoriales cilíndricos, entre los que los haces radiales, azimutales y espirales son casos particulares, y corresponden a nuestro caso de estudio.


Como ya tenemos claro, el campo transmitido en el plano focal es proporcional a 

\begin{equation}
  \mathscr F \{ \mathbb T \mathbb E_0 \}
\end{equation}

al $\mathbb E_0$ no depender de $\phi$ tenemos que el campo es proporcional a 

\begin{equation}
  \int_0^{\infty} \mathscr F_\phi \{ \mathbb T\} \mathbb E_0 r'dr'
\end{equation}



Para calcular $\mathscr{F}_\phi$ se usa la identidad 


\begin{equation}
\int_0^{2\pi} e^{im\theta}e^{-ir\cos(\theta - \theta')}\,d\theta,
=
2\pi(-i)^m e^{im\theta'} J_m(r)
\qquad (m\in\mathbb Z)
\end{equation}

que es la expansión de Jacobi-Anger con un cambio de variable \cite{Watson1944},


de la cual se obtiene que 



\begin{equation}
  \mathscr{F}_\phi\{\mathbb{T}\}(q, r', \varphi)
  =
  \pi\,J_0(qr')\,t_+\,\mathbb{I}_\perp
  -
  \pi\,J_2(qr')\,t_-\,\mathbb{V}(\varphi),
  \label{eq:Fphi_M_circular}
\end{equation}

donde

\begin{equation}
\mathbb V(\varphi)
=
\begin{bmatrix}
0 & e^{-2i\varphi}
\\
e^{2i\varphi} & 0
\end{bmatrix}.
\label{eq:def_V_mixing}
\end{equation}


Quedando ahora las funciones de Bessel $J_0$ y $J_2$ bajo una integral respecto a la coordenada radial con peso la primera potencia de dicha coordenada, el campo transmitido al plano focal se escribe en términos de las transformadas de Hankel de orden 0 y 2


\begin{equation}
\boxed{
\begin{aligned}
E'_L(\vb r,f)
&=
\frac{\pi n}{i\lambda f z_0}
e^{inkf}
e^{i\frac{nk}{2f}r^2}
\Big[
\mathscr H_0\{t_+\mathscr E_L\}(q)
-
e^{-2i\varphi}
\mathscr H_2\{t_-\mathscr E_R\}(q)
\Big],
\\
E'_R(\vb r,f)
&=
\frac{\pi n}{i\lambda f z_0}
e^{inkf}
e^{i\frac{nk}{2f}r^2}
\Big[
-e^{2i\varphi} \mathscr H_2\{t_-\mathscr E_L\}(q)
+
\mathscr H_0\{t_+\mathscr E_R\}(q)
\Big].
\end{aligned}
\label{eq:E_L_E_R_Hankel}
}
\end{equation}

O de forma equivalente, en notación matricial,

\begin{equation}
\boxed{
\begin{pmatrix} E'_L \\ E'_R \end{pmatrix}
=
\frac{\pi n}{i\lambda f z_0}\,e^{inkf}\,e^{i\frac{nk}{2f}r^2}
\left[
\mathscr H_0[t_+]\,\mathbb I_\perp
-
\mathscr H_2[t_-]\,\mathbb V(\varphi)
\right]
\begin{pmatrix} \mathscr E_L \\ \mathscr E_R \end{pmatrix}.
}
\label{eq:E_L_E_R_Hankel_matrix}
\end{equation}

donde $\mathscr H_m[f]$ denota el operador que multiplica por $f(r')$ y aplica la transformada de Hankel de orden $m$, y $t_\pm$ están definidas en \eqref{eq:t_plus_t_minus_def}.

Los resultados para la pupila circular uniforme en las distintas polarizaciones se desarrollan en el apéndice~\ref{apx:focal_forms}.

El término proporcional a $\mathscr{H}_0[t_+]\,\mathbb{I}_\perp$ es diagonal: transforma $\ket{L}\to\ket{L}$ y $\ket{R}\to\ket{R}$ sin mezclar helicidades, reproduciendo el resultado de la difracción escalar modulado por $t_+$. La aparición de la aberración vectorial queda asociada de forma explícita al término de mezcla proporcional a $t_-=t_p-t_s$, que acopla $\ket{L}$ con $\ket{R}$ mediante factores de fase azimutal $e^{\pm 2i\varphi}$. Si $t_p=t_s$, dicho acoplamiento desaparece.

Las series paraxiales del apéndice~\ref{apx:ovoids_series} muestran que, al orden cuadrático en la coordenada pupilar,

\begin{equation}
t_+(\rho')
\approx
\frac{4n_0}{n_0+n_i}
+
\frac{
n_0(n_0-n_i)(z_0-z_i)^2
}{
z_0^2z_i^2(n_0+n_i)^2
}\,\rho'^2,
\qquad
t_-(\rho')
\approx
\frac{
n_0(n_0-n_i)^2(z_0-z_i)^2
}{
z_0^2z_i^2(n_0+n_i)^3
}\,\rho'^2,
\label{eq:t_plus_t_minus_paraxial_second_order}
\end{equation}

Así, $t_+(0)=4n_0/(n_0+n_i)\neq 0$ mientras que $t_-(0)=0$: el coeficiente de mezcla es nulo sobre el eje y crece cuadráticamente alejándose de él. A esto se añade que $J_2(0)=0$, de modo que tanto el coeficiente $t_-$ como la función de Bessel $J_2$ se anulan simultáneamente en el origen, suprimiendo la conversión de helicidad en el eje. En contraste, $J_0(0)=1$ y $t_+(0)\neq 0$ garantizan que el término diagonal sí contribuye en el eje, generando un brillo central en el plano focal.

Esto se confirma para una entrada circular en la figura~\ref{fig:circular-plano-focal-componentes}. En el plano focal domina el canal de la misma helicidad, con centro brillante, mientras que en la componente de helicidad opuesta aparece la conversión parcial inducida por $t_-$: un anillo con centro oscuro, consecuencia directa de que $t_-(0)=0$ y $J_2(0)=0$ suprimen la conversión sobre el eje. El mapa de polarización de la figura~\ref{fig:circular-plano-focal-mapa-polarizacion} muestra cómo esta conversión modifica espacialmente el estado local de polarización, y la figura~\ref{fig:circular-plano-focal-mapas-escalares} separa dicha variación en los ángulos $\chi$ y $\psi$.

\begin{figure}[H]
  \centering
  \thesisincludegraphics[width=\textwidth]
    {figures/results_of_simulation/vecdiff_final/circular/circular_componentes.png}
    {figures/fast/circular_componentes.jpg}
  \caption{Componentes circulares $|E_L|$, $|E_R|$ e intensidad para incidencia circular derecha: campo sobre la pupila (fila superior) y campo en el plano focal (fila inferior). Dioptrio estigmático con $n_0=1$, $n_i=1.5$, $z_0=-10\,\mathrm{mm}$, $z_i=6\,\mathrm{mm}$, $\lambda=532\,\mathrm{nm}$ y pupila uniforme. La componente $E_L$ en el plano focal exhibe la conversión parcial de helicidad inducida por $t_-$, con estructura anular y centro oscuro.}
  \label{fig:circular-plano-focal-componentes}
\end{figure}

\begin{figure}[H]
  \centering
  \thesisincludegraphics[width=\textwidth]
    {figures/results_of_simulation/vecdiff_final/circular/circular_mapa_polarizacion.png}
    {figures/fast/circular_mapa_polarizacion.jpg}
  \caption{Mapas de polarización para incidencia circular: campo sobre la pupila (izquierda) y campo en el plano focal (derecha). El fondo muestra $|E_x|^2+|E_y|^2$ y las elipses el estado local de polarización.}
  \label{fig:circular-plano-focal-mapa-polarizacion}
\end{figure}

\begin{figure}[H]
  \centering
  \thesisincludegraphics[width=0.9\textwidth]
    {figures/results_of_simulation/vecdiff_final/circular/circular_angulos_elipse.png}
    {figures/fast/circular_angulos_elipse.jpg}
  \caption{Ángulos de la elipse de polarización para incidencia circular: elipticidad $\chi$ (izquierda) y orientación del eje mayor $\psi$ (derecha), sobre la pupila (fila superior) y en el plano focal (fila inferior). Las regiones en blanco corresponden a zonas de intensidad demasiado baja, donde estos parámetros dejan de estar bien definidos.}
  \label{fig:circular-plano-focal-mapas-escalares}
\end{figure}


% ==== Representación Cilíndrica === 



% ···········································································
\subsection{Campo eléctrico en el plano focal en la representación cilíndrica}
% ···········································································

La representación cilíndrica es la base natural asociada a las direcciones locales de polarización radial y azimutal. En esta base, la acción transversal de los coeficientes de Fresnel es diagonal antes de efectuar la propagación difractiva. Es decir, si el campo incidente sobre la superficie de referencia se escribe como

\begin{equation}
  \mathbb E_0(\rho',\phi')
  =
  \mathscr E_r(\rho',\phi')\,\uvec r(\phi')
  +
  \mathscr E_\phi(\rho',\phi')\,\uvec \phi(\phi'),
\end{equation}

entonces el campo transmitido transversalmente queda dado por

\begin{equation}
  \vb E_t(\rho',\phi')
  =
  t_p(\rho')\mathscr E_r(\rho',\phi')\,\uvec r(\phi')
  +
  t_s(\rho')\mathscr E_\phi(\rho',\phi')\,\uvec \phi(\phi').
  \label{eq:T_cylindrical_basis}
\end{equation}

En notación operatorial,

\begin{equation}
  \mathbb T
  =
  t_p\,\ket{\uvec r}\bra{\uvec r}
  +
  t_s\,\ket{\uvec\phi}\bra{\uvec\phi},
  \label{eq:M_in_cylindrical_basis}
\end{equation}

donde

\begin{equation}
  \uvec r(\phi')
  =
  \cos\phi'\,\uvec x
  +
  \sin\phi'\,\uvec y,
  \qquad
  \uvec\phi(\phi')
  =
  -\sin\phi'\,\uvec x
  +
  \cos\phi'\,\uvec y.
  \label{eq:cylindrical_unit_vectors}
\end{equation}

Por tanto, reemplazando en la expresión general del campo en el plano focal,

\begin{equation}
\boxed{
  \vb E'_\perp(\vb r,f)
  =
  \frac{n}{i\lambda f z_0}
  e^{inkf}
  e^{i\frac{nk}{2f}r^2}
  \mathscr F
  \left\{
    t_p\mathscr E_r\,\uvec r
    +
    t_s\mathscr E_\phi\,\uvec\phi
  \right\}
  \left(
    \frac{n}{\lambda f}\vb r
  \right).
}
\label{eq:E_prime_cylindrical_general}
\end{equation}

Este resultado es todav\'ia general: no se ha supuesto que las componentes escalares \(\mathscr E_r\) y \(\mathscr E_\phi\) sean axialmente sim\'etricas.

\subsubsection{Reducción del problema a componentes axialmente simétricas}

Supongamos ahora que las componentes cil\'indricas del campo incidente no dependen de \(\phi'\), es decir,

\begin{equation}
  \mathscr E_r=\mathscr E_r(\rho'),
  \qquad
  \mathscr E_\phi=\mathscr E_\phi(\rho').
\end{equation}

En coordenadas polares de frecuencia, la transformada de Fourier bidimensional se escribe como

\begin{equation}
  \mathscr F\{f\}(q,\varphi)
  =
  \int_0^\infty
  \int_0^{2\pi}
  f(\rho',\phi')
  e^{-iq\rho'\cos(\phi'-\varphi)}
  d\phi'\,\rho' d\rho',
  \label{eq:fourier_polar_cylindrical}
\end{equation}

con

\begin{equation}
  q=\frac{n}{\lambda f}r.
\end{equation}

Para efectuar la integral angular usamos la identidad de Jacobi--Anger

\begin{equation}
  \int_0^{2\pi}
  e^{im\phi'}
  e^{-iq\rho'\cos(\phi'-\varphi)}
  d\phi'
  =
  2\pi(-i)^m e^{im\varphi}J_m(q\rho'),
  \qquad m\in\mathbb Z.
  \label{eq:jacobi_anger_cylindrical}
\end{equation}

De ella se deducen, para \(m=1\),

\begin{equation}
  \int_0^{2\pi}
  \cos\phi'\,
  e^{-iq\rho'\cos(\phi'-\varphi)}
  d\phi'
  =
  -2\pi i\cos\varphi\,J_1(q\rho'),
\end{equation}

\begin{equation}
  \int_0^{2\pi}
  \sin\phi'\,
  e^{-iq\rho'\cos(\phi'-\varphi)}
  d\phi'
  =
  -2\pi i\sin\varphi\,J_1(q\rho').
\end{equation}

Por lo tanto,

\begin{equation}
  \mathscr F_\phi\{\uvec r(\phi')\}
  =
  -2\pi i\,J_1(q\rho')\,\uvec r(\varphi),
  \label{eq:Fphi_rhat}
\end{equation}

y, de manera an\'aloga,

\begin{equation}
  \mathscr F_\phi\{\uvec\phi(\phi')\}
  =
  -2\pi i\,J_1(q\rho')\,\uvec\phi(\varphi).
  \label{eq:Fphi_phihat}
\end{equation}

N\'otese que lo que se transforma de esta forma no es el operador dy\'adico \(\mathbb T\) por s\'i solo, sino el campo transmitido \(\mathbb T\mathbb E_0\). Sustituyendo \eqref{eq:Fphi_rhat} y \eqref{eq:Fphi_phihat} en \eqref{eq:E_prime_cylindrical_general}, se obtiene

\begin{equation}
\boxed{
  \vb E'_\perp(\vb r,f)
  =
  -\frac{2\pi n}{\lambda f z_0}
  e^{inkf}
  e^{i\frac{nk}{2f}r^2}
  \left[
    \mathscr H_1\{t_p\mathscr E_r\}(q)\,\uvec r(\varphi)
    +
    \mathscr H_1\{t_s\mathscr E_\phi\}(q)\,\uvec\phi(\varphi)
  \right].
}
\label{eq:E_field_cylindrical_Hankel}
\end{equation}

Aqu\'i

\begin{equation}
  \mathscr H_m\{g\}(q)
  =
  \int_0^\infty
  g(\rho')J_m(q\rho')\,\rho' d\rho'
  \label{eq:Hankel_definition_cylindrical}
\end{equation}

denota la transformada de Hankel de orden \(m\).

Equivalentemente, si escribimos el resultado en la base cil\'indrica del plano focal,

\begin{equation}
  \vb E'_\perp
  =
  E'_r\,\uvec r(\varphi)
  +
  E'_\phi\,\uvec\phi(\varphi),
\end{equation}

entonces

\begin{equation}
\boxed{
\begin{pmatrix}
E'_r\\
E'_\phi
\end{pmatrix}
=
-\frac{2\pi n}{\lambda f z_0}
e^{inkf}
e^{i\frac{nk}{2f}r^2}
\begin{pmatrix}
\mathscr H_1[t_p] & 0\\
0 & \mathscr H_1[t_s]
\end{pmatrix}
\begin{pmatrix}
\mathscr E_r\\
\mathscr E_\phi
\end{pmatrix}.
}
\label{eq:E_cylindrical_Hankel_matrix}
\end{equation}

Los resultados para la pupila circular uniforme se desarrollan en el apéndice~\ref{apx:focal_forms}.


Es claro cómo en esta representación la transformación que sufre el campo al refractarse en la interfaz $\Sigma$ es diagonal, esto es que los modos radiales y azimutales de la polarización son modos propios de este proceso de refracción a través de una superficie. La modulación radial está dada por la transformada de Hankel de orden 1 de los coeficientes de Fresnel; dado que la transformada de Hankel de orden uno es una integral ponderada por la función de Bessel $J_1$, no hay contribución central del campo en el eje óptico, como se observa en las figuras~\ref{fig:radial-plano-focal-componentes} y~\ref{fig:azimutal-plano-focal-componentes}. Esto es consecuencia de la homogeneidad del campo en la polarización radial y azimutal.

Podemos decir que \textit{un campo radial o azimutal simétrico que sea axialmente simétrico, en su propagación hasta el plano focal, da lugar a un campo sin contribución central}.


\begin{figure}[H]
  \centering
  \thesisincludegraphics[width=\textwidth]
    {figures/results_of_simulation/vecdiff_final/cilindrico/radial_componentes.png}
    {figures/fast/radial_componentes.jpg}
  \caption{Componentes cilíndricas $|E_r|$, $|E_\phi|$ e intensidad para incidencia radial: campo sobre la pupila (fila superior) y campo en el plano focal (fila inferior). Dioptrio estigmático con $n_0=1$, $n_i=1.5$, $z_0=-10\,\mathrm{mm}$, $z_i=6\,\mathrm{mm}$, $\lambda=532\,\mathrm{nm}$ y pupila uniforme. La transmisión es diagonal en esta base: no aparece componente azimutal, y la estructura de Hankel de orden uno produce el centro oscuro.}
  \label{fig:radial-plano-focal-componentes}
\end{figure}

\begin{figure}[H]
  \centering
  \thesisincludegraphics[width=\textwidth]
    {figures/results_of_simulation/vecdiff_final/cilindrico/radial_mapa_polarizacion.png}
    {figures/fast/radial_mapa_polarizacion.jpg}
  \caption{Mapas de polarización para incidencia radial: campo sobre la pupila (izquierda) y campo en el plano focal (derecha). El fondo muestra $|E_x|^2+|E_y|^2$ y los glifos el estado local de polarización, que permanece radial punto a punto.}
  \label{fig:radial-plano-focal-mapa-polarizacion}
\end{figure}

\begin{figure}[H]
  \centering
  \thesisincludegraphics[width=\textwidth]
    {figures/results_of_simulation/vecdiff_final/cilindrico/azimutal_componentes.png}
    {figures/fast/azimutal_componentes.jpg}
  \caption{Componentes cilíndricas $|E_r|$, $|E_\phi|$ e intensidad para incidencia azimutal: campo sobre la pupila (fila superior) y campo en el plano focal (fila inferior), con los mismos parámetros de la figura~\ref{fig:radial-plano-focal-componentes}. La polarización azimutal es modo propio de la transmisión (coeficiente $t_s$) y el plano focal exhibe de nuevo el anillo con centro oscuro.}
  \label{fig:azimutal-plano-focal-componentes}
\end{figure}

\begin{figure}[H]
  \centering
  \thesisincludegraphics[width=\textwidth]
    {figures/results_of_simulation/vecdiff_final/cilindrico/azimutal_mapa_polarizacion.png}
    {figures/fast/azimutal_mapa_polarizacion.jpg}
  \caption{Mapas de polarización para incidencia azimutal: campo sobre la pupila (izquierda) y campo en el plano focal (derecha). El estado local permanece azimutal punto a punto.}
  \label{fig:azimutal-plano-focal-mapa-polarizacion}
\end{figure}




% ==== REPRESENTACIÓN CARTESIANA === 


% ···········································································
\subsection{Campo eléctrico en el plano focal en la representación cartesiana}
\label{subsec:representacion_cartesiana}
% ···········································································

Consideremos ahora el campo incidente escrito en la base cartesiana transversal,

\begin{equation}
  \mathscr E_\perp
  =
  \mathscr E_x\,\hat{\mathbf x}
  +
  \mathscr E_y\,\hat{\mathbf y}.
\end{equation}

Las direcciones locales asociadas a la base de Fresnel sobre la pupila son

\begin{equation}
\hat{\mathbf r}(\phi')
=
\cos\phi'\,\hat{\mathbf x}
+
\sin\phi'\,\hat{\mathbf y},
\qquad
\hat{\boldsymbol\phi}(\phi')
=
-\sin\phi'\,\hat{\mathbf x}
+
\cos\phi'\,\hat{\mathbf y}.
\label{eq:cartesian_polar_unit_vectors}
\end{equation}

Por tanto, las componentes del campo incidente en la base local
$\{\hat{\mathbf r},\hat{\boldsymbol\phi}\}$ son

\begin{equation}
\begin{aligned}
  \mathscr E_r
  &=
  \mathscr E_x\cos\phi'
  +
  \mathscr E_y\sin\phi',
  \\
  \mathscr E_\phi
  &=
  -
  \mathscr E_x\sin\phi'
  +
  \mathscr E_y\cos\phi'.
\end{aligned}
\label{eq:cartesian_to_local_fresnel_components}
\end{equation}

Después de la transmisión en la superficie dióptrica, el campo transversal local queda dado por

\begin{equation}
  \mathbf T
  =
  t_p\mathscr E_r\,\hat{\mathbf r}
  +
  t_s\mathscr E_\phi\,\hat{\boldsymbol\phi}.
\label{eq:T_cartesian_local_fresnel}
\end{equation}

Sustituyendo \eqref{eq:cartesian_to_local_fresnel_components} en
\eqref{eq:T_cartesian_local_fresnel}, y proyectando de nuevo sobre la base
cartesiana, se obtiene

\begin{equation}
\begin{aligned}
T_x
&=
\left(
t_p\cos^2\phi'
+
t_s\sin^2\phi'
\right)\mathscr E_x
+
\left(
t_p-t_s
\right)
\sin\phi'\cos\phi'\,\mathscr E_y,
\\
T_y
&=
\left(
t_p-t_s
\right)
\sin\phi'\cos\phi'\,\mathscr E_x
+
\left(
t_p\sin^2\phi'
+
t_s\cos^2\phi'
\right)\mathscr E_y.
\end{aligned}
\label{eq:Tx_Ty_before_trig_reduction}
\end{equation}

Así,

\begin{equation}
  \begin{pmatrix}
    T_x\\
    T_y
  \end{pmatrix}
  =
  \mathbb T_{xy}
  \begin{pmatrix}
    \mathscr E_x\\
    \mathscr E_y
  \end{pmatrix}.
\label{eq:T_cartesian_matrix_def}
\end{equation}

Usando las identidades

\begin{equation}
\cos^2\phi'
=
\frac{1}{2}
\left(
1+\cos2\phi'
\right),
\qquad
\sin^2\phi'
=
\frac{1}{2}
\left(
1-\cos2\phi'
\right),
\qquad
\sin\phi'\cos\phi'
=
\frac{1}{2}\sin2\phi',
\end{equation}

la matriz de transmisión en la base cartesiana adopta la forma

\begin{equation}
\boxed{
  \mathbb T_{xy}
  =
  \frac{t_p+t_s}{2}\,\mathbb I_\perp
  +
  \frac{t_p-t_s}{2}
  \begin{pmatrix}
    \cos 2\phi' & \sin 2\phi' \\
    \sin 2\phi' & -\cos 2\phi'
  \end{pmatrix}.
}
\label{eq:M_in_cartesian_basis}
\end{equation}

Usando las combinaciones $t_+$ y $t_-$ definidas en \eqref{eq:t_plus_t_minus_def}, y definiendo

\begin{equation}
  \mathbb U(\phi')
  =
  \begin{pmatrix}
    \cos 2\phi' & \sin 2\phi' \\
    \sin 2\phi' & -\cos 2\phi'
  \end{pmatrix},
\label{eq:U_phi_cartesian_def}
\end{equation}

podemos escribir simplemente

\begin{equation}
  \mathbb T_{xy}
  =
  \frac{1}{2}t_+\,\mathbb I_\perp
  +
  \frac{1}{2}t_-\,\mathbb U(\phi').
\label{eq:M_cartesian_tplus_tminus}
\end{equation}

La descomposición anterior separa el operador de transmisión en una parte
isotrópica, proporcional a $\mathbb I_\perp$, y una parte anisotrópica de
orden angular dos, proporcional a $\mathbb U(\phi')$.

El campo transmitido en el plano focal puede escribirse entonces como

\begin{equation}
\boxed{
\mathbf E'_\perp(\mathbf r,f)
=
\frac{n e^{inkf}}{i\lambda f z_0}
e^{i\frac{nk}{2f}r^2}
\mathscr F
\left[
\mathbb T_{xy}(\rho',\phi')
\begin{pmatrix}
\mathscr E_x\\
\mathscr E_y
\end{pmatrix}
\right]
\left(
\frac{n}{\lambda f}\mathbf r
\right).
}
\label{eq:E_prime_cartesian_general_Fourier}
\end{equation}

De manera equivalente, separando la integración angular de la integración radial,

\begin{equation}
\mathbf E'_\perp(\mathbf r,f)
=
\frac{n e^{inkf}}{i\lambda f z_0}
e^{i\frac{nk}{2f}r^2}
\int_0^\infty
\mathscr F_\phi\{\mathbb T_{xy}\}
(q,\rho',\varphi)
\begin{pmatrix}
\mathscr E_x(\rho',\phi')\\
\mathscr E_y(\rho',\phi')
\end{pmatrix}
\rho'\,d\rho',
\label{eq:E_prime_cartesian_general_radial_integral}
\end{equation}

donde

\begin{equation}
  q
  =
  \frac{n r}{\lambda f}.
\label{eq:q_cartesian_def}
\end{equation}

\subsubsection{Reducción del problema a componentes axialmente simétricas}

Supongamos ahora que las componentes cartesianas del campo incidente son
axialmente simétricas, es decir,

\begin{equation}
  \mathscr E_x
  =
  \mathscr E_x(\rho'),
  \qquad
  \mathscr E_y
  =
  \mathscr E_y(\rho').
\label{eq:cartesian_axial_symmetry_assumption}
\end{equation}

En este caso, toda la dependencia angular de la integral se encuentra en
$\mathbb T_{xy}$. Para calcular la transformada angular se usan

\begin{equation}
  \cos2\phi'
  =
  \frac{1}{2}
  \left(
  e^{2i\phi'}
  +
  e^{-2i\phi'}
  \right),
  \qquad
  \sin2\phi'
  =
  \frac{1}{2i}
  \left(
  e^{2i\phi'}
  -
  e^{-2i\phi'}
  \right),
\end{equation}

junto con la identidad de Jacobi--Anger,

\begin{equation}
\int_0^{2\pi}
e^{im\phi'}
e^{-iq\rho'\cos(\phi'-\varphi)}
\,d\phi'
=
2\pi(-i)^m e^{im\varphi}J_m(q\rho'),
\qquad
m\in\mathbb Z.
\label{eq:jacobi_anger_cartesian_use}
\end{equation}

Puesto que $(-i)^0=1$ y $(-i)^{\pm2}=-1$, se obtiene

\begin{equation}
\boxed{
  \mathscr F_\phi\{\mathbb T_{xy}\}
  (q,\rho',\varphi)
  =
  \pi J_0(q\rho')\,t_+(\rho')\,\mathbb I_\perp
  -
  \pi J_2(q\rho')\,t_-(\rho')\,\mathbb W(\varphi),
}
\label{eq:Fphi_M_cartesian}
\end{equation}

donde hemos definido el operador de mezcla cartesiano

\begin{equation}
\boxed{
  \mathbb W(\varphi)
  =
  \begin{pmatrix}
    \cos 2\varphi & \sin 2\varphi \\
    \sin 2\varphi & -\cos 2\varphi
  \end{pmatrix}.
}
\label{eq:def_W_mixing}
\end{equation}

La sustitución de \eqref{eq:Fphi_M_cartesian} en
\eqref{eq:E_prime_cartesian_general_radial_integral} da

\begin{equation}
\boxed{
\mathbf E'_\perp(\mathbf r,f)
=
\frac{\pi n e^{inkf}}{i\lambda f z_0}
e^{i\frac{nk}{2f}r^2}
\left[
\mathbb I_\perp
\mathscr H_0\{t_+\mathscr E_\perp\}(q)
-
\mathbb W(\varphi)
\mathscr H_2\{t_-\mathscr E_\perp\}(q)
\right].
}
\label{eq:E_field_cartesian_azimuthal}
\end{equation}

Aquí la transformada de Hankel actúa componente a componente, es decir,

\begin{equation}
\mathscr H_m\{t_\pm\mathscr E_\perp\}
=
\begin{pmatrix}
\mathscr H_m\{t_\pm\mathscr E_x\}\\
\mathscr H_m\{t_\pm\mathscr E_y\}
\end{pmatrix}.
\label{eq:Hankel_componentwise_cartesian}
\end{equation}

Por componentes, \eqref{eq:E_field_cartesian_azimuthal} se escribe como

\begin{equation}
\boxed{
\begin{aligned}
E'_x(\mathbf r,f)
=
\frac{\pi n e^{inkf}}{i\lambda f z_0}
e^{i\frac{nk}{2f}r^2}
\Big[
&
\mathscr H_0\{t_+\mathscr E_x\}(q)
\\
&
-
\cos2\varphi\,
\mathscr H_2\{t_-\mathscr E_x\}(q)
\\
&
-
\sin2\varphi\,
\mathscr H_2\{t_-\mathscr E_y\}(q)
\Big],
\\[1em]
E'_y(\mathbf r,f)
=
\frac{\pi n e^{inkf}}{i\lambda f z_0}
e^{i\frac{nk}{2f}r^2}
\Big[
&
\mathscr H_0\{t_+\mathscr E_y\}(q)
\\
&
-
\sin2\varphi\,
\mathscr H_2\{t_-\mathscr E_x\}(q)
\\
&
+
\cos2\varphi\,
\mathscr H_2\{t_-\mathscr E_y\}(q)
\Big].
\end{aligned}
}
\label{eq:Ex_Ey_Hankel_cartesian}
\end{equation}

O bien, en notación matricial,


\begin{equation}
\boxed{
\begin{pmatrix}
E'_x\\
E'_y
\end{pmatrix}
=
\frac{\pi n e^{inkf}}{i\lambda f z_0}
e^{i\frac{nk}{2f}r^2}
\begin{pmatrix}
\mathscr H_0[t_+] - \cos2\varphi\,\mathscr H_2[t_-]
&
-\sin2\varphi\,\mathscr H_2[t_-]
\\
-\sin2\varphi\,\mathscr H_2[t_-]
&
\mathscr H_0[t_+] + \cos2\varphi\,\mathscr H_2[t_-]
\end{pmatrix}
\begin{pmatrix}
\mathscr E_x\\
\mathscr E_y
\end{pmatrix}.
}
\label{eq:Ex_Ey_Hankel_cartesian_matrix}
\end{equation}


\begin{equation}
\begin{pmatrix}
E'_x\\
E'_y
\end{pmatrix}
=
\frac{\pi n e^{inkf}}{i\lambda f z_0}
e^{i\frac{nk}{2f}r^2}
\left[
\mathscr H_0[t_+] I
-
\mathscr H_2[t_-]\mathbb W(\varphi) \right]
\begin{pmatrix}
\mathscr E_x\\
\mathscr E_y
\end{pmatrix}.
\label{eq:Ex_Ey_Hankel_cartesian_matrix_separated}
\end{equation}




En esta expresión, $\mathscr H_m[t_\pm]$ denota el operador que primero
multiplica por $t_\pm(\rho')$ y luego aplica la transformada de Hankel de
orden $m$.

Los resultados para la pupila homogénea y los casos particulares de polarización se desarrollan en el apéndice~\ref{apx:focal_forms}.

El término $\mathscr{H}_0[t_+]\,\mathbb{I}_\perp$ es isótropo: preserva la dirección de polarización ($\hat{\mathbf x}\to\hat{\mathbf x}$, $\hat{\mathbf y}\to\hat{\mathbf y}$) y reproduce el resultado escalar modulado por $t_+$. La aparición de la aberración vectorial se identifica en el término anisótropo proporcional a $t_-=t_p-t_s$, que introduce la mezcla angular de orden dos a través de $\mathbb{W}(\varphi)$ y de las contribuciones $\mathscr{H}_2[t_-]$. Si $t_p=t_s$, sobrevive solo la parte isotrópica y el término vectorial aberrante desaparece.

Al igual que en la representación circular, la supresión axial de la mezcla queda garantizada por la aproximación paraxial \eqref{eq:t_plus_t_minus_paraxial_second_order}: $t_-(0)=0$ y $J_2(0)=0$ anulan simultáneamente la contribución de $\mathscr{H}_2[t_-]$ sobre el eje, mientras que $t_+(0)\neq 0$ y $J_0(0)=1$ aseguran un brillo central en el canal isótropo.

La consecuencia más visible de $\mathbb{W}(\varphi)$ es la aparición de una componente cruzada en el plano focal con dependencia azimutal de orden dos. Para una entrada puramente $\hat{\mathbf x}$, el término $-\sin 2\varphi\,\mathscr{H}_2[t_-\mathscr{E}_x]$ genera una componente $E_y'$ con patrón de cuatro lóbulos, varios órdenes de magnitud menor que $E_x'$, como se observa en la figura~\ref{fig:lineal-plano-focal-componentes}. El mapa de la figura~\ref{fig:lineal-plano-focal-mapa-polarizacion} evidencia la rotación local inducida por esta componente cruzada, y la figura~\ref{fig:lineal-plano-focal-mapas-escalares} separa dicha variación en los ángulos $\chi$ y $\psi$: la elipticidad permanece nula en todo el plano focal y la aberración vectorial se manifiesta exclusivamente como una rotación local de la orientación, hecho que se justificó en la sección~\ref{subsec:polarizacion_plano_focal}.

\begin{figure}[H]
  \centering
  \thesisincludegraphics[width=\textwidth]
    {figures/results_of_simulation/vecdiff_final/lineal/lineal_componentes.png}
    {figures/fast/lineal_componentes.jpg}
  \caption{Componentes cartesianas $|E_x|$, $|E_y|$ e intensidad para incidencia lineal $\hat{\mathbf x}$: campo sobre la pupila (fila superior) y campo en el plano focal (fila inferior). Dioptrio estigmático con $n_0=1$, $n_i=1.5$, $z_0=-10\,\mathrm{mm}$, $z_i=6\,\mathrm{mm}$, $\lambda=532\,\mathrm{nm}$ y pupila uniforme. La componente $|E_y'|$ corresponde a la polarización cruzada, con el patrón de cuatro lóbulos $\sin 2\varphi$ predicho por \eqref{eq:Ex_Ey_Hankel_cartesian}.}
  \label{fig:lineal-plano-focal-componentes}
\end{figure}

\begin{figure}[H]
  \centering
  \thesisincludegraphics[width=\textwidth]
    {figures/results_of_simulation/vecdiff_final/lineal/lineal_mapa_polarizacion.png}
    {figures/fast/lineal_mapa_polarizacion.jpg}
  \caption{Mapas de polarización para incidencia lineal $\hat{\mathbf x}$: campo sobre la pupila (izquierda) y campo en el plano focal (derecha). El fondo muestra $|E_x|^2+|E_y|^2$ y las flechas la orientación local de la polarización, evidenciando la rotación inducida por la componente cruzada.}
  \label{fig:lineal-plano-focal-mapa-polarizacion}
\end{figure}

\begin{figure}[H]
  \centering
  \thesisincludegraphics[width=0.9\textwidth]
    {figures/results_of_simulation/vecdiff_final/lineal/lineal_angulos_elipse.png}
    {figures/fast/lineal_angulos_elipse.jpg}
  \caption{Ángulos de la elipse de polarización para incidencia lineal $\hat{\mathbf x}$: elipticidad $\chi$ (izquierda) y orientación del eje mayor $\psi$ (derecha), sobre la pupila (fila superior) y en el plano focal (fila inferior). La elipticidad es nula en todo el plano focal: el estado permanece lineal y la aberración vectorial se reduce a una rotación local de la orientación. Las regiones en blanco corresponden a zonas de intensidad demasiado baja, donde estos parámetros dejan de estar bien definidos.}
  \label{fig:lineal-plano-focal-mapas-escalares}
\end{figure}

\subsubsection{Efecto del régimen vectorial sobre el tamaño del punto focal}
\label{subsubsec:spot_vectorial}

El resultado \eqref{eq:Ex_Ey_Hankel_cartesian} permite cuantificar de forma
cerrada cómo el régimen vectorial modifica el tamaño del punto focal. Para
incidencia lineal $\hat{\mathbf x}$ con amplitud axialmente simétrica
($\mathscr E_y=0$, $\mathscr E_x=\mathscr E_x(\rho')$), y escribiendo por
brevedad $H_0=\mathscr H_0\{t_+\mathscr E_x\}(q)$ y
$H_2=\mathscr H_2\{t_-\mathscr E_x\}(q)$, la intensidad transversal en el
plano focal resulta

\begin{equation}
  I_\perp(\vb r, f)
  \propto
  |H_0|^2
  +
  |H_2|^2
  -
  2\operatorname{Re}(H_0 H_2^*)\cos 2\varphi.
\label{eq:I_spot_linear_deformation}
\end{equation}

Si la polarización incidente forma un ángulo $\psi_0$ con el eje $x$, el
patrón es el mismo con $\varphi\to\varphi-\psi_0$: el punto focal rota
rígidamente con la polarización, sin grados de libertad angulares adicionales.
La estructura de \eqref{eq:I_spot_linear_deformation} implica que el lóbulo
central deja de ser circular: sus ejes principales son las direcciones paralela
y perpendicular a la polarización, con cortes $I_\parallel\propto|H_0-H_2|^2$ e
$I_\perp^{(90°)}\propto|H_0+H_2|^2$, y el radio único del caso escalar se
reemplaza por dos radios principales $R_\parallel$ y $R_\perp$, medidos a mitad
de la altura del máximo central. Cerca del eje se tiene
$\operatorname{Re}(H_0H_2^*)>0$, de modo que el término de interferencia
sustrae intensidad del flanco del lóbulo a lo largo de la polarización y la
agrega en la dirección ortogonal: respecto de la referencia escalar
—$|H_0|^2$, con radio $R_{\mathrm{esc}}$— el punto se contrae a lo largo de la
polarización y se expande en la dirección perpendicular,

\begin{equation}
  R_\parallel < R_{\mathrm{esc}} < R_\perp.
\label{eq:spot_radii_ordering}
\end{equation}

La figura~\ref{fig:spot-cortes-principales} muestra los cortes por los ejes
principales para el dioptrio de vidrio usado en los resultados anteriores, a
apertura alta: la contracción es de $\sim 2.6\,\%$ y la expansión de
$\sim 2.8\,\%$, con una elipticidad del contorno $e=R_\perp/R_\parallel\approx
1.06$. La deformación crece con la apertura y con el contraste de índice, pues
ambos incrementan el canal $H_2$. La figura~\ref{fig:spot-radios-anotados}
presenta el caso de un dioptrio rápido de alto índice ($n_i=2.4$,
$z_0=-2\,\mathrm{mm}$, $z_i=6\,\mathrm{mm}$, semiapertura objeto
$\approx 57°$): los radios medidos son $R_\parallel=0.245\lambda$ y
$R_\perp=0.279\lambda$ frente a $R_{\mathrm{esc}}=0.260\lambda$, es decir,
$e\approx 1.14$. Para incidencia circular la intensidad es
$I\propto|H_0|^2+|H_2|^2$, el promedio azimutal del caso lineal: el punto
permanece circular y solo se ensancha levemente ($R=0.261\lambda$ en el mismo
dioptrio), pues el canal cruzado $|H_2|^2$ es un anillo nulo en el eje. Cabe
señalar que estos resultados corresponden a la intensidad transversal; a
aperturas muy altas la contribución longitudinal $I_z$, determinada por la
condición de transversalidad, modifica adicionalmente la forma del lóbulo
total.

\begin{figure}[H]
  \centering
  \thesisincludegraphics[width=0.75\textwidth]
    {figures/results_of_simulation/vecdiff_final/spot/spot_cortes_principales.png}
    {figures/fast/spot_cortes_principales.jpg}
  \caption{Cortes de la intensidad focal por los ejes principales para
  incidencia lineal en el dioptrio de vidrio ($n_0=1$, $n_i=1.5$,
  $z_0=-10\,\mathrm{mm}$, $z_i=6\,\mathrm{mm}$, $\lambda=532\,\mathrm{nm}$,
  pupila uniforme a apertura alta). El corte paralelo a la polarización se
  contrae y el perpendicular se expande respecto de la referencia escalar
  apodizada ($t_-=0$); la referencia escalar ideal ($t_p=t_s=1$) muestra que la
  apodización de Fresnel $t_+(\rho')$ también ensancha el punto.}
  \label{fig:spot-cortes-principales}
\end{figure}

\begin{figure}[H]
  \centering
  \thesisincludegraphics[width=\textwidth]
    {figures/results_of_simulation/vecdiff_final/spot/spot_radios_anotados.png}
    {figures/fast/spot_radios_anotados.jpg}
  \caption{Radios principales del lóbulo central, medidos a mitad de altura,
  para un dioptrio rápido de alto índice ($n_i=2.4$, $z_0=-2\,\mathrm{mm}$,
  $z_i=6\,\mathrm{mm}$, semiapertura objeto $\approx 57°$). Izquierda:
  incidencia lineal con polarización a $20°$; el contorno medido (blanco) es
  elíptico, con $R_\parallel=0.245\lambda$, $R_\perp=0.279\lambda$
  ($e\approx1.14$) frente al radio escalar $R_{\mathrm{esc}}=0.260\lambda$
  (cian, a trazos), y los ejes principales siguen a la polarización. Derecha:
  incidencia circular; el punto permanece circular con $R=0.261\lambda$,
  levemente mayor que el escalar.}
  \label{fig:spot-radios-anotados}
\end{figure}
